% AI Finance 1 — 1일차 강의 슬라이드 (Beamer, metropolis)
% 빌드: xelatex AI_Finance_1.tex  (두 번)
\documentclass[aspectratio=169,10pt]{beamer}
% 공통 스타일 (metropolis, 한글, 색, 표, 그림)
\usetheme[progressbar=frametitle,numbering=fraction,block=fill]{metropolis}

\usepackage{fontspec}
\usepackage[hangul]{kotex}
\setmainfont{Apple SD Gothic Neo}
\setsansfont{Apple SD Gothic Neo}[BoldFont={* Bold}]
\setmonofont{Menlo}[Scale=0.88]
\setmainhangulfont{Apple SD Gothic Neo}
\setsanshangulfont{Apple SD Gothic Neo}[BoldFont={* Bold}]
\setmonohangulfont{Apple SD Gothic Neo}

\usepackage{booktabs,tabularx,array,colortbl}
\usepackage{pgfplots}
\pgfplotsset{compat=1.18,/pgf/number format/1000 sep={\mathord{,}}}
\usepackage{tikz}
\usetikzlibrary{arrows.meta,positioning,shapes.misc}

% 색 (metropolis 기본과 맞춤)
\definecolor{ink}{HTML}{23373B}
\definecolor{accent}{HTML}{EB811B}
\definecolor{teal}{HTML}{14B03D}
\definecolor{blue}{HTML}{2B5EA7}
\definecolor{soft}{HTML}{F2F4F8}
\definecolor{muted}{HTML}{6B7785}
\setbeamercolor{alerted text}{fg=accent}
\setbeamercolor{block title}{bg=ink!12,fg=ink}
\setbeamercolor{block body}{bg=soft}

\newcommand{\src}[1]{{\scriptsize\color{muted}#1}}
\newcommand{\num}[1]{{\color{accent}\textbf{#1}}}
\newcommand{\tagbox}[2]{{\setlength{\fboxsep}{2.5pt}\colorbox{#1}{\textcolor{white}{\scriptsize\bfseries #2}}}}
\renewcommand{\arraystretch}{1.25}
\tikzset{card/.style={rectangle,rounded corners=4pt,draw=#1,line width=0.9pt,fill=#1!6,
               minimum width=6.6cm,minimum height=2.25cm,align=left,text width=6.2cm,inner sep=6pt},
  box/.style={rounded corners=3pt,draw=#1,fill=#1!8,line width=0.8pt,minimum width=4.2cm,minimum height=0.9cm,align=center,font=\small},
  arr/.style={-{Stealth[length=5pt]},line width=0.8pt,muted}}

% Claude 요청문·코드: \begin{Verbatim} … \end{Verbatim} (frame은 [fragile])
\usepackage{fvextra}
\fvset{breaklines=true,breakanywhere=true,breaksymbolleft={},fontsize=\scriptsize,frame=leftline,framerule=2pt,rulecolor=\color{accent},xleftmargin=4pt,framesep=6pt}


\title{SEC 공시에서 내 DB까지}
\subtitle{AI Finance 1 — 받고, 바꾸고, 넣고, 묻는다}
\date{2026년 11월 28일 · 대우관 415호}

\begin{document}

\maketitle

% ─────────────────────────────────────────────
\begin{frame}{오늘의 흐름}
\centering
\begin{tikzpicture}[node distance=0.55cm]
\node[box=accent,minimum width=3.0cm,text width=2.8cm] (a) {\textbf{1. 받기}\\\scriptsize SEC 10월 공시 원문};
\node[box=blue,minimum width=3.0cm,text width=2.8cm,right=of a] (b) {\textbf{2. 바꾸기}\\\scriptsize HTML → 텍스트};
\node[box=teal,minimum width=3.0cm,text width=2.8cm,right=of b] (c) {\textbf{3. 넣기}\\\scriptsize 내 Azure DB};
\node[box=ink,minimum width=3.0cm,text width=2.8cm,right=of c] (d) {\textbf{4. 묻기}\\\scriptsize SQL · 임베딩};
\draw[arr] (a) -- (b); \draw[arr] (b) -- (c); \draw[arr] (c) -- (d);
\end{tikzpicture}
\vspace{0.8em}

\footnotesize\renewcommand{\arraystretch}{1.1}
\begin{tabularx}{0.92\textwidth}{@{}l r X l@{}}
\toprule
\textbf{시간} & \textbf{분} & \textbf{활동} & \textbf{기다리는 동안} \\
\midrule
15:00 & 5 & 시작 · 사전 준비 확인 & \\
15:05 & 25 & \tagbox{accent}{실습 1} SEC에서 10월 공시 받기 & 다운로드 약 20분 \\
15:30 & 25 & \tagbox{blue}{실습 2} HTML을 텍스트로 & \\
15:55 & 10 & 휴식 & \\
16:05 & 40 & \tagbox{teal}{실습 3} Azure에 내 DB 만들기 & 서버 생성 약 6분 \\
16:45 & 55 & \tagbox{ink}{실습 4} 임베딩 · SQL · 질의 & 벡터화 약 16분 \\
17:40 & 20 & 정리 · 퀴즈 (오픈북) & \\
\bottomrule
\end{tabularx}
\end{frame}

% ─────────────────────────────────────────────
\begin{frame}{Claude와 일하는 규칙}
\begin{columns}[T]
\begin{column}{0.48\textwidth}
\begin{block}{Claude는 코드를 만들고, 설명하고, 고친다}
각 단계는 \textbf{정해진 코드}로 돈다.\\
Claude Code는 명령을 실행하기 전에 \textbf{허락을 묻는다} — 무엇을 하려는지 읽고 허락한다.
\end{block}
\end{column}
\begin{column}{0.48\textwidth}
\begin{alertblock}{올리지 않는 것}
수업은 \textbf{공개된 SEC 공시}만 다룬다.\\
회사 자료·고객 정보는 Claude나 Azure에 올리지 않는다.
\end{alertblock}
\end{column}
\end{columns}
\vspace{0.8em}
\begin{alertblock}{\texttt{.env}는 비밀 파일}
이메일, 이후 DB 비밀번호와 키가 들어간다. 채팅에 붙여 넣거나 Claude에게 열어 보라고 하지 않는다.\\
\src{SEC는 자동으로 공시를 받는 프로그램에 연락처를 밝히라고 요구한다. 이메일은 SEC에 보내는 요청의 머리글에만 쓰인다.}
\end{alertblock}
\end{frame}

% ─────────────────────────────────────────────
\section{사전 준비 확인}

\begin{frame}{사전 준비 — 손들어 확인}
\footnotesize\renewcommand{\arraystretch}{1.12}
\begin{tabularx}{\textwidth}{@{}r X X@{}}
\toprule
\textbf{\#} & \textbf{할 일} & \textbf{확인 방법} \\
\midrule
1 & \textbf{Azure 무료 계정} (11월 5일 이후 가입 — 크레딧 30일) & portal.azure.com에서 \textbf{구독}이 보인다 \\
2 & \textbf{VS Code} + \textbf{Claude Code} 확장 최신 & Claude Code가 ``안녕''에 답한다 \\
3 & \textbf{PostgreSQL} 확장 \texttt{ms-ossdata.vscode-pgsql} \src{(게시자 Microsoft)} & 왼쪽에 코끼리 아이콘 \\
4 & 작업 폴더 \texttt{문서/jpm-class} — 같은 창에서 파일 → 폴더 열기 & 탐색기에 \texttt{JPM-CLASS} \\
5 & \textbf{uv} (요청 1) & \texttt{Python 3.12.x} \\
6 & \textbf{Azure CLI} (요청 2) & \texttt{az version} \\
7 & 터미널에서 \texttt{az login} & 구독 이름이 나온다 \\
8 & \textbf{실습 키트} → \texttt{jpm-class} 안에 풀기 & \texttt{jpm\_kit/README.md} \\
9 & \textbf{\texttt{.env}}: \texttt{SEC\_CONTACT} 한 줄 (영문 이름 + 이메일) & \texttt{.env} 파일과 그 한 줄 \\
\bottomrule
\end{tabularx}
\vspace{0.3em}
\src{준비가 안 된 사람은 옆 사람과 짝을 지어 시작합니다.}
\end{frame}

\begin{frame}[fragile]{설치는 Claude에게 — 요청 1, 2}
\textbf{요청 1} \src{uv (Python 실행 도구)}
\begin{Verbatim}
Python 패키지 관리자 uv 가 설치되어 있는지 확인하고, 없으면 공식 설치 방법으로 설치해줘.
설치 후 `uv run --python 3.12 python --version` 을 실행해 결과를 보여줘.
\end{Verbatim}
\vspace{0.4em}
\textbf{요청 2} \src{Azure CLI}
\begin{Verbatim}
Azure CLI 가 설치되어 있는지 확인하고, 없으면 이 운영체제에 맞는 공식 방법(Windows: winget, Mac: Homebrew)으로 설치해줘.
설치 후 `az version` 결과를 보여줘. 로그인(az login)은 내가 직접 할게.
\end{Verbatim}
\end{frame}

\begin{frame}[fragile]{막혔을 때}
\begin{columns}[T]
\begin{column}{0.47\textwidth}
\begin{block}{Windows}
\small
``앱이 디바이스를 변경하도록 허용하시겠습니까?'' → \textbf{예}.\\[3pt]
Claude가 \texttt{uv}·\texttt{az}를 ``찾을 수 없다''고 하면 VS Code를 \textbf{완전히 닫았다 다시 열고} 같은 요청을 다시.\\[3pt]
문서가 \textbf{OneDrive}에 있으면 \texttt{C:\textbackslash Users\textbackslash <내 이름>\textbackslash jpm-class}에 만든다.
\end{block}
\end{column}
\begin{column}{0.49\textwidth}
\begin{block}{Mac에 Homebrew가 없으면}
\small
Claude는 관리자 비밀번호를 입력할 수 없다. 터미널에서 \textbf{직접}:
\end{block}
\vspace{-0.4em}
\begin{Verbatim}[fontsize=\tiny]
/bin/bash -c "$(curl -fsSL https://raw.githubusercontent.com/Homebrew/install/HEAD/install.sh)"
\end{Verbatim}
\small
끝에 나오는 \textbf{Next steps}의 \texttt{echo} 두 줄, \texttt{eval} 한 줄을 그대로 붙여 넣기 → \texttt{brew -{}-version} → VS Code \textbf{⌘Q} 후 요청 2를 다시.
\end{column}
\end{columns}
\vspace{0.6em}
\begin{alertblock}{\texttt{.env}는 VS Code 탐색기에서}
\small
\texttt{.env.example} 복사 → 붙여넣기 → 이름 바꾸기 \texttt{.env}. 예: \texttt{SEC\_CONTACT=Gildong Hong gildong@example.com} — \textbf{이름은 영문}, 따옴표 없이. \src{메모장은 \texttt{.env.txt}로 저장하는 일이 많다.}
\end{alertblock}
\end{frame}

% ─────────────────────────────────────────────
% 1일차 · 실습 1 — SEC에서 10월 공시 받기 (실습_SEC_10월_수집.md)
\section{실습 1. SEC에서 10월 공시 받기}

% ─────────────────────────────────────────────
\begin{frame}{목표: 공시 원문 한 달 치를 SEC에서 직접}
\begin{block}{이번 실습}
JPMorgan이 \num{2026년 10월}에 낸 \textbf{424B2 공시 원문}을 \textbf{SEC에서 직접} 받는 코드를 Claude와 만들고,\\
\textbf{내가 읽고, 내가 실행하고,} 받은 것이 맞는지 \textbf{확인}합니다.
\end{block}
\vspace{0.6em}
\centering\small 이 데이터가 이후 모든 실습의 출발점입니다.\\[6pt]
\begin{tikzpicture}[node distance=0.35cm]
\node[box=accent,minimum width=2.3cm] (a) {SEC 원문};
\node[box=ink,minimum width=2.3cm,right=of a] (b) {텍스트 변환};
\node[box=ink,minimum width=2.3cm,right=of b] (c) {규칙 추출};
\node[box=ink,minimum width=2.3cm,right=of c] (d) {상품 표};
\node[box=ink,minimum width=2.3cm,right=of d] (e) {검색};
\draw[arr] (a) -- (b); \draw[arr] (b) -- (c); \draw[arr] (c) -- (d); \draw[arr] (d) -- (e);
\end{tikzpicture}
\end{frame}

% ─────────────────────────────────────────────
\begin{frame}{할 일과 시간}
\small
\renewcommand{\arraystretch}{1.1}
\begin{tabularx}{\textwidth}{@{}X l l@{}}
\toprule
\textbf{할 일} & \textbf{어디서} & \textbf{시간} \\
\midrule
1. 연락처가 \texttt{.env}에 있는지 확인 (사전 준비 9) & VS Code 탐색기 & 1분 \\
2. Claude에게 코드 만들게 하기 & Claude 패널 & 10분 \\
3. 코드를 읽고 요구사항 확인 & Claude 패널 & 10분 \\
4. 목록만 받기, 원문 3건만 받기 (시험) & Claude 패널 & 5분 \\
5. 10월 전체 받기 (백그라운드, 약 20분) & Claude 패널 & 기다리는 동안 다음 설명 \\
6. 받은 것 검증 & Claude 패널 & 5분 \\
\bottomrule
\end{tabularx}
\vspace{0.5em}
\begin{alertblock}{SEC 공정 접근 규칙}
한 IP에서 \num{초당 10건}까지. 넘으면 SEC가 그 IP를 10분 이상 막습니다(\texttt{403}).\\
이 강의실은 노트북마다 IP가 달라 각자 받을 수 있지만, 내 코드는 \num{초당 2건}으로 여유 있게 받습니다.
\end{alertblock}
\end{frame}

% ─────────────────────────────────────────────
\begin{frame}[fragile]{1. 연락처가 \texttt{.env}에 있는지 확인 \src{\normalsize 사전 준비 9에서 적어 옴}}
\small
\noindent SEC는 자동으로 받는 프로그램에 \textbf{연락처(이름과 이메일)}를 밝히라고 요구합니다. 없으면 \texttt{403}으로 거부합니다.
\vspace{0.4em}
\begin{columns}[T]
\begin{column}{0.50\textwidth}
\textbf{\texttt{문서/jpm-class/jpm\_kit/.env}}
\begin{Verbatim}
SEC_CONTACT=Gildong Hong gildong@example.com
\end{Verbatim}
\begin{itemize}\setlength\itemsep{1pt}
\item VS Code 탐색기에서 열어 \textbf{있는지만} 확인하고 탭을 닫습니다 (없으면 지금 적습니다)
\item \texttt{.env.txt}로 보이면 이름을 \texttt{.env}로 바꿉니다
\end{itemize}
\end{column}
\begin{column}{0.46\textwidth}
\begin{alertblock}{한글 이름은 안 됩니다}
따옴표 없이 \texttt{영문 이름 빈칸 이메일} 그대로.\\
요청 머리글(User-Agent)에는 영문·숫자·기호만 쓸 수 있어 오류로 멈춥니다.\\
\src{\texttt{UnicodeEncodeError: 'latin-1' codec can't encode}}
\end{alertblock}
\begin{block}{이메일은 채팅에 쓰지 않습니다}
코드가 \texttt{.env}를 직접 읽습니다. Claude는 \texttt{.env}를 열지 않습니다 \src{(키트의 \texttt{CLAUDE.md} 규칙)}.
\end{block}
\end{column}
\end{columns}
\end{frame}

% ─────────────────────────────────────────────
\begin{frame}[fragile]{2. Claude에게 코드 만들게 하기}
\small
\noindent VS Code에서 폴더 \texttt{jpm-class}를 연 상태로, 왼쪽 \textbf{Claude Code} 패널에 아래 요구사항을 \textbf{그대로} 붙여 넣습니다.
\vspace{0.3em}
\begin{Verbatim}
jpm_kit/sec_oct.py 를 만들어줘. 아직 실행하지 마.
목적: JPMorgan이 2026-10-01 ~ 2026-10-31에 SEC에 제출한 424B2 공시 원문을 받는다.

요구사항
1. 대상 목록: https://data.sec.gov/submissions/CIK0000019617.json (JPMorgan Chase & Co.)과
   CIK0001665650.json (발행 자회사)의 filings.recent 에서 form 이 424B2 이고 filingDate 가 10월인 것.
   같은 공시가 두 회사에 함께 등록되므로 accessionNumber 로 합친다.
2. 원문 주소: https://www.sec.gov/Archives/edgar/data/19617/<accessionNumber에서 '-'를 뺀 것>/<primaryDocument>
3. 연락처: jpm_kit/.env 의 SEC_CONTACT 값을 User-Agent 로 쓴다. 값을 화면에 출력하거나 다른 파일에 쓰지 않는다.
   값이 비었거나 영문·숫자·기호(ASCII) 밖의 글자가 있으면 요청을 보내기 전에 멈추고 "SEC_CONTACT를 영문 이름과 이메일로 적으세요"라고만 출력한다.
   너(Claude)는 .env 를 열지 않는다. 코드가 실행될 때 읽는다.
\end{Verbatim}
\end{frame}

\begin{frame}[fragile]{2. Claude에게 코드 만들게 하기 (계속)}
\begin{Verbatim}
4. 속도: 요청은 한 번에 하나씩, 초당 2건 이하.
5. 403 또는 429 를 받으면 즉시 전체를 멈추고 이유를 출력한다. 다른 방법으로 다시 시도하지 않는다.
6. 저장: jpm_kit/sec_oct/html/<accessionNumber>.html. 이미 있는 파일은 다시 받지 않는다(중단 후 이어 받기).
7. 기록: jpm_kit/sec_oct/filings.csv (accession, filingDate, 등록 회사, 원문 주소)와
   jpm_kit/sec_oct/download_log.csv (accession, 결과, 바이트 수, sha256, 받은 시각).
8. 옵션: --list-only (목록만 만들기), --limit N (원문 N건만 받기).
9. 100건마다 진행 상황(받은 수, 경과 시간, 남은 시간 추정)을 한 줄로 출력한다.
10. Python 표준 라이브러리만 쓰고, 실행은 `uv run jpm_kit/sec_oct.py` 로 한다.
다 만들면 코드가 위 10가지를 각각 어디서 지키는지 줄 번호로 알려줘.
\end{Verbatim}
\end{frame}

% ─────────────────────────────────────────────
\begin{frame}{3. 코드를 읽고 요구사항 확인}
\small
\noindent Claude가 알려 준 줄 번호를 \textbf{직접 열어} 확인하고 표를 채웁니다. 하나라도 \alert{``아니요''}면 Claude에게 고치게 합니다.
\vspace{0.3em}

\begin{tabularx}{\textwidth}{@{}c X p{2.2cm} p{2.2cm}@{}}
\toprule
\textbf{\#} & \textbf{요구사항} & \textbf{지키나?} & \textbf{줄 번호} \\
\midrule
1 & 두 회사 목록을 합친다 \src{(한 회사만 보면 몇 건이 빠질까?)} & & \\
3 & 연락처를 출력하지 않는다 & & \\
4 & 초당 2건 이하, 한 번에 하나 & & \\
5 & \texttt{403}·\texttt{429}면 즉시 멈춘다 & & \\
6 & 이미 받은 파일은 건너뛴다 & & \\
\bottomrule
\end{tabularx}
\vspace{0.4em}
\begin{block}{확인 질문}
\textit{``이 코드가 SEC에 보내는 요청을 하나씩 순서대로 말해줘. 10월 전체를 받으면 요청이 모두 몇 번이고 몇 분 걸려?''}
\end{block}
\end{frame}

% ─────────────────────────────────────────────
\begin{frame}[fragile]{4. 시험: 목록만, 그다음 3건만}
\begin{columns}[T]
\begin{column}{0.40\textwidth}
\begin{Verbatim}
uv run jpm_kit/sec_oct.py --list-only
uv run jpm_kit/sec_oct.py --limit 3
\end{Verbatim}
\vspace{0.6em}
\begin{alertblock}{403이 나오면}
\textbf{그 자리에서 멈추고 손을 듭니다.}\\
다시 실행하지 않습니다.
\end{alertblock}
\end{column}
\begin{column}{0.56\textwidth}
\small
\begin{tabularx}{\textwidth}{@{}X p{2.4cm}@{}}
\toprule
\textbf{확인} & \textbf{내 값} \\
\midrule
10월 424B2 공시 수 & \\[4pt]
그중 두 회사에 함께 등록된 수 & \\[4pt]
받은 3건의 accession & \\[4pt]
\texttt{html} 폴더에서 1건을 브라우저로 열었을 때 보이는 제목 & \\[4pt]
\bottomrule
\end{tabularx}
\end{column}
\end{columns}
\end{frame}

% ─────────────────────────────────────────────
\begin{frame}[fragile]{5. 10월 전체 받기}
\small
\noindent Claude Code 패널에서:
\begin{Verbatim}
jpm_kit/sec_oct.py 를 백그라운드로 실행하고, 진행 상황을 주기적으로 알려줘.
\end{Verbatim}
\vspace{0.3em}
\begin{columns}[T]
\begin{column}{0.48\textwidth}
\begin{block}{얼마나 걸리나}
초당 2건이면 \num{약 1,900건}에 \num{20분 남짓}.\\
노트북 전원을 연결하고 덮지 않습니다.\\
기다리는 동안 다음 실습(HTML을 텍스트로 바꾸기)의 설명을 듣습니다.
\end{block}
\end{column}
\begin{column}{0.48\textwidth}
\begin{block}{멈추거나 실패하면}
중간에 멈추면 같은 명령을 다시 실행합니다(이어 받기).\\
끝났을 때 \texttt{failed}가 0이 아니면 한 번 더 실행합니다. 받은 파일은 건너뛰고 실패한 것만 다시 받습니다.
\end{block}
\end{column}
\end{columns}
\end{frame}

% ─────────────────────────────────────────────
\begin{frame}[fragile]{5. 403이 나오거나 시간 안에 못 받으면}
\begin{alertblock}{강사가 올린 압축 파일로 대신합니다}
\href{https://www.dropbox.com/scl/fi/a7njti2aj7789ccrj0cyw/10-_.zip?rlkey=tk8t7s50x2vsotvejaeh8vfit&dl=1}{\texttt{10월\_원문.zip}} \src{(51MB)}을 \texttt{jpm\_kit/sec\_oct/} 폴더에 풉니다.\\[2pt]
\src{받은 파일 이름이 \texttt{10-\_.zip}처럼 보일 수 있음}
\end{alertblock}
\vspace{0.6em}
\centering
\begin{tikzpicture}[node distance=0.6cm]
\node[box=accent,minimum width=3.2cm] (z) {\texttt{10월\_원문.zip}};
\node[box=ink,minimum width=3.6cm,right=1.2cm of z,yshift=0.6cm] (f) {\texttt{filings.csv}};
\node[box=ink,minimum width=3.6cm,right=1.2cm of z,yshift=-0.6cm] (h) {\texttt{html/} \src{(1,859건)}};
\node[box=teal,minimum width=2.8cm,right=1.2cm of f,yshift=-0.6cm] (v) {검증(6)은 똑같이};
\draw[arr] (z.east) -- (f.west); \draw[arr] (z.east) -- (h.west);
\draw[arr] (f.east) -- (v.west); \draw[arr] (h.east) -- (v.west);
\end{tikzpicture}
\end{frame}

% ─────────────────────────────────────────────
\begin{frame}[fragile]{6. 받은 것이 맞는지 검증}
\begin{Verbatim}
uv run jpm_kit/check_oct.py
\end{Verbatim}
\small
\begin{columns}[T]
\begin{column}{0.44\textwidth}
\noindent 키트의 검증 스크립트가 내 \texttt{sec\_oct/}를 \textbf{강사가 같은 날 받은 목록과 원문}에 대조합니다.\\[4pt]
\src{SEC는 원문을 보낼 때마다 끝에 봇 차단용 \texttt{<script>} 한 줄을 바꿔 붙이므로, 그 줄만 빼고 비교합니다.}
\end{column}
\begin{column}{0.52\textwidth}
\begin{tabularx}{\textwidth}{@{}X p{2.2cm}@{}}
\toprule
\textbf{확인} & \textbf{내 값} \\
\midrule
목록: 강사 목록과 같은 accession 수 / 빠진 수 / 더 있는 수 & \\[6pt]
원문: 같은 파일 수 / 다른 파일 수 / 없는 파일 수 & \\[6pt]
\bottomrule
\end{tabularx}
\end{column}
\end{columns}
\end{frame}

% ─────────────────────────────────────────────
\begin{frame}{생각해 보기}
\begin{enumerate}\setlength\itemsep{10pt}
\item 발행 자회사(CIK 1665650) 목록만 썼다면 \alert{몇 건이 빠졌을까요?} 왜 빠질까요?
\item 같은 공시를 오늘 받고 내일 받으면 파일의 sha256이 다릅니다.\\
그런데도 \alert{``같은 원문''}이라고 말할 수 있는 근거는?
\item 받은 1,900여 건은 모두 \alert{서로 다른 상품}일까요? \src{(다음 실습에서 확인합니다)}
\end{enumerate}
\end{frame}

% 1일차 · 실습 2. HTML을 텍스트로 바꾸기 (원본: 실습_HTML_텍스트_변환.md)
\section{실습 2. HTML을 텍스트로 바꾸기}

% ─────────────────────────────────────────────
\begin{frame}{변환기가 하는 일 — 오늘의 목표}
\begin{columns}[T]
\begin{column}{0.46\textwidth}
\begin{block}{목표}
받은 공시 원문(HTML)을 텍스트로 바꾸고, 변환기가 \textbf{무엇을 바꾸고 무엇을 바꾸지 않는지} 이해합니다.
\end{block}
\begin{alertblock}{고치지 않고 읽고 확인}
변환기 \texttt{jpm\_kit/htmltext.py}는 강사가 만들어 2026년 1–10월 공시 \num{19,285건}으로 검사한 것을 그대로 씁니다.
\end{alertblock}
\end{column}
\begin{column}{0.50\textwidth}
\small
\begin{tabularx}{\textwidth}{@{}X r@{}}
\toprule
\textbf{할 일} & \textbf{시간} \\
\midrule
1. 왜 바꾸나, HTML 두 종류 \src{(설명)} & 5분 \\
2. 변환기가 하는 일 \src{(설명)} & 10분 \\
3. 내 10월 원문 전체를 바꾸고 검사 & 5분 \\
4. 두 종류를 하나씩 원문 화면과 나란히 보기 & 10분 \\
5. Claude에게 변환기 코드 한 부분 설명 듣기 & 5분 \\
\bottomrule
\end{tabularx}
\end{column}
\end{columns}
\end{frame}

% ─────────────────────────────────────────────
\begin{frame}{왜 텍스트로 바꾸나}
다음 단계들(규칙으로 값 뽑기, 상품 표, 검색, AI에게 읽히기)은 모두 \textbf{글자}를 다룹니다.\\
HTML에는 글자 말고도 \alert{태그, 글꼴, 위치 정보}가 섞여 있습니다.
\vspace{0.8em}

\small
\begin{tabularx}{\textwidth}{@{}l X X@{}}
\toprule
 & \textbf{HTML} & \textbf{텍스트} \\
\midrule
크기 (중앙값) & 일반 HTML \num{119KB}, PDF 변환형 \num{743KB} & \textcolor{accent}{\textbf{50–60KB}} \\
\texttt{Pricing Date: October 16, 2025} & 표의 두 칸에 나뉘어 태그 사이에 흩어짐 & \texttt{Pricing Date: \textbar{} October 16, 2025} 한 줄 \\
\texttt{Linked to} & \texttt{Linke\textless{}font …\textgreater{}d to \textless{}/font\textgreater{}}\newline\src{(낱말 가운데 태그)} & \texttt{Linked to} \\
\bottomrule
\end{tabularx}
\end{frame}

% ─────────────────────────────────────────────
\begin{frame}{공시 HTML은 두 종류입니다}
\small 원본이 그렇게 생겼습니다. 공시 대행사가 어떻게 만들었느냐에 따라 갈립니다.
\vspace{0.4em}

\centering
\begin{tikzpicture}
% ① 일반 HTML
\node[card=teal,minimum height=4.8cm,text width=6.2cm,anchor=north] (a) at (0,0) {};
\node[anchor=north west,font=\bfseries] at ([xshift=4pt,yshift=-4pt]a.north west) {① 일반 HTML \textcolor{muted}{\small(문서 구조가 있음)}};
\node[anchor=north west,font=\scriptsize,text width=6cm] at ([xshift=4pt,yshift=-22pt]a.north west) {처음부터 HTML로 작성. 문단은 \texttt{<p>}, 표는 \texttt{<table><tr><td>}.};
\node[draw=teal,fill=white,rounded corners=2pt,minimum width=5.6cm,minimum height=0.45cm,font=\scriptsize] (p1) at ([yshift=-2.05cm]a.north) {\texttt{<p>} 문단};
\node[draw=teal,fill=white,rounded corners=2pt,minimum width=5.6cm,minimum height=0.8cm,anchor=north] (t1) at ([yshift=-0.15cm]p1.south) {};
\node[draw=teal,fill=teal!15,minimum width=2.6cm,minimum height=0.5cm,font=\tiny,anchor=west] at ([xshift=0.15cm]t1.west) {\texttt{<td>}Pricing Date:};
\node[draw=teal,fill=teal!15,minimum width=2.6cm,minimum height=0.5cm,font=\tiny,anchor=east] at ([xshift=-0.15cm]t1.east) {\texttt{<td>}October 16, 2025};
\node[font=\tiny,text=teal,anchor=south east] at ([yshift=-2pt]t1.north east) {\texttt{<tr>} 한 행};
\node[draw=teal,fill=white,rounded corners=2pt,minimum width=5.6cm,minimum height=0.45cm,font=\scriptsize,anchor=north] at ([yshift=-0.15cm]t1.south) {\texttt{<p>} 문단};
\node[anchor=south,font=\scriptsize,text=teal] at ([yshift=4pt]a.south) {태그가 \textbf{문단·표·칸}을 알려 준다};
% ② PDF 변환형
\node[card=accent,minimum height=4.8cm,text width=6.2cm,anchor=north] (b) at (7.2,0) {};
\node[anchor=north west,font=\bfseries] at ([xshift=4pt,yshift=-4pt]b.north west) {② PDF를 HTML로 바꾼 형식};
\node[anchor=north west,font=\scriptsize,text width=6cm] at ([xshift=4pt,yshift=-22pt]b.north west) {Word·조판 프로그램으로 PDF를 만든 뒤 변환 도구로 HTML로. 문단·표 태그가 없음.};
\begin{scope}[shift={([yshift=-2.0cm]b.north)}]
\foreach \y/\w/\x in {0/4.6/-0.3, -0.45/3.8/-0.7, -0.9/4.4/-0.4, -1.35/1.6/-1.8, -1.35/1.4/1.2}{
  \node[draw=accent,fill=accent!15,minimum width=\w cm,minimum height=0.3cm,inner sep=0pt] at (\x,\y) {};
}
\draw[muted,densely dotted] (-2.75,0.25) -- (-2.75,-1.6);
\node[font=\tiny,text=muted,anchor=east] at (-2.75,0) {\texttt{bottom}};
\node[font=\tiny,text=muted,anchor=north] at (-0.3,-1.55) {\texttt{left:935.99px; bottom:12842.0px}};
\end{scope}
\node[anchor=south,font=\scriptsize,text=accent] at ([yshift=4pt]b.south) {화면의 \textbf{한 줄 = 상자 하나}, \textbf{좌표}로 놓임};
\end{tikzpicture}
\end{frame}

% ─────────────────────────────────────────────
\begin{frame}[fragile]{두 종류의 원문 — 실제 모습}
\begin{block}{① 일반 HTML}
\begin{Verbatim}
<tr><td>Pricing Date:</td><td>October 16, 2025</td></tr>
\end{Verbatim}
\end{block}
\begin{block}{② PDF를 HTML로 바꾼 형식}
\begin{Verbatim}
<div class="t …" style="left:935.99px; bottom:12842.0px; font-size:379.2px">Uncapped Digital Barrier Notes Linke<font style="letter-spacing:0.364px">d to </font>the Lesser </div>
<div class="t …" style="left:935.99px; bottom:12405.2px; font-size:379.2px">Performing of …</div>
\end{Verbatim}
\end{block}
\small 브라우저로 열면 PDF와 똑같이 보이지만, \alert{``이 글자들이 한 문장''이라는 정보가 파일에 없습니다}. 좌표를 보고 다시 맞춰야 합니다.

\vspace{0.3em}
\src{2026년 1–10월 공시 19,285건에서 ②는 \textbf{5,896건(31\%)}이었고, 모두 대행사 한 곳(접수번호 앞자리 \texttt{0001213900})에서 나왔습니다.}
\end{frame}

% ─────────────────────────────────────────────
\begin{frame}{원칙은 하나}
\begin{center}
\large \textbf{글자는 하나도 바꾸지 않고, 글자 사이에 무엇을 넣을지만 정한다}
\end{center}
\vspace{0.2em}
\centering
\begin{tikzpicture}
\node[box=ink,minimum width=3.2cm] (g) at (5.4,1.5) {글자 \quad ? \quad 글자};
\node[box=teal,minimum width=2.6cm] (a) at (0,0) {붙임};
\node[box=teal,minimum width=2.6cm] (b) at (3.6,0) {공백};
\node[box=blue,minimum width=2.6cm] (c) at (7.2,0) {줄바꿈};
\node[box=accent,minimum width=2.6cm] (d) at (10.8,0) {칸 구분 \texttt{\textvisiblespace|\textvisiblespace}};
\foreach \n in {a,b,c,d} \draw[arr] (g.south) -- (\n.north);
\end{tikzpicture}
\vspace{0.6em}

\raggedright
\begin{alertblock}{태그를 지우는 것만으로는 이 결정을 할 수 없습니다}
태그가 있던 자리에 공백을 넣으면 \texttt{Linke d to}가 되고,
아무것도 안 넣으면 \texttt{Pricing Date:October 16}이 됩니다.
\end{alertblock}
\end{frame}

% ─────────────────────────────────────────────
\begin{frame}{① 일반 HTML \texttt{(dom\_to\_text)} — 1/2}
\small
\begin{tabularx}{\textwidth}{@{}>{\raggedright\arraybackslash}X >{\raggedright\arraybackslash}p{5.0cm} >{\raggedright\arraybackslash}p{3.0cm}@{}}
\toprule
\textbf{규칙} & \textbf{원문} & \textbf{결과} \\
\midrule
문단·목록 같은 덩어리(\texttt{<p>}, \texttt{<div>}, \texttt{<li>} …)는 줄을 바꾼다 & \texttt{<p>…notes.</p><p>The notes…</p>} & 두 줄 \\
표의 한 행은 한 줄, 칸 사이는 \texttt{\textvisiblespace|\textvisiblespace} & \texttt{<tr><td>CUSIP:</td><td>48136HP69</td></tr>} & \texttt{CUSIP: \textbar{} 48136HP69} \\
글자 가운데의 태그는 그냥 붙인다 & \texttt{Linke<font>d to</font>} & \texttt{Linked to} \\
\texttt{<br>}(줄바꿈)이 문장·날짜 한가운데면 공백, 아니면 줄바꿈 & \texttt{June 29,<br>2028} & \texttt{June 29, 2028} \\
\bottomrule
\end{tabularx}
\end{frame}

\begin{frame}{① 일반 HTML \texttt{(dom\_to\_text)} — 2/2}
\small
\begin{tabularx}{\textwidth}{@{}>{\raggedright\arraybackslash}X >{\raggedright\arraybackslash}p{5.0cm} >{\raggedright\arraybackslash}p{3.0cm}@{}}
\toprule
\textbf{규칙} & \textbf{원문} & \textbf{결과} \\
\midrule
HTML 소스 안의 줄바꿈은 화면에서 공백일 뿐 & \texttt{the\textbackslash{}nnotes} & \texttt{the notes} \\
화면에 안 보이는 것은 뺀다 & \texttt{<script>}, \texttt{<style>}, \texttt{<!-{}- 주석 -{}->} & (없음) \\
폭을 지정한 넓은 빈 상자는 다른 항목 사이의 빈칸 & \texttt{Investments —<font style="width:400pt"> </font>PS-1} & \texttt{Investments — \textbar{} PS-1} \\
\bottomrule
\end{tabularx}
\end{frame}

% ─────────────────────────────────────────────
\begin{frame}{② PDF를 HTML로 바꾼 형식 \texttt{(pdf\_to\_text)} — 1/2}
\small
\begin{tabularx}{\textwidth}{@{}>{\raggedright\arraybackslash}p{3.6cm} >{\raggedright\arraybackslash}X >{\raggedright\arraybackslash}p{4.4cm}@{}}
\toprule
\textbf{규칙} & \textbf{판단 근거} & \textbf{결과 예} \\
\midrule
높이(\texttt{bottom})가 같은 상자들은 한 행 & 좌표 & 표의 한 행 \\
같은 행에서 멀리 떨어진 상자는 다른 칸 & 간격 > 글자 크기의 \num{1.5배} & \texttt{Per note \textbar{} \$1,000} \\
상자 안의 빈 간격 조각(\texttt{\_})은 폭으로 판단 & 글자 크기의 \num{15\%} 이하 = 붙임, 그보다 넓으면 공백, \num{5배} 넘으면 칸 & \texttt{January , 2026 \textbar{} Registration Statement Nos.} \\
\bottomrule
\end{tabularx}
\vspace{0.8em}

\centering
\begin{tikzpicture}[font=\scriptsize]
\draw[muted,line width=0.6pt] (0,0) -- (12,0);
\fill[teal!60] (0,-0.12) rectangle (1.8,0.12);
\fill[blue!60] (1.8,-0.12) rectangle (4.2,0.12);
\fill[accent!70] (4.2,-0.12) rectangle (12,0.12);
\node[below] at (0,-0.15) {0};
\node[below] at (1.8,-0.15) {0.15 × 글자 크기};
\node[below] at (4.2,-0.15) {5 × 글자 크기};
\node[above] at (0.9,0.15) {붙임};
\node[above] at (3.0,0.15) {공백};
\node[above] at (8.1,0.15) {칸 \texttt{\textbar{}}};
\node[anchor=west,text=muted] at (12.1,0) {빈 조각 폭};
\end{tikzpicture}
\end{frame}

\begin{frame}{② PDF를 HTML로 바꾼 형식 \texttt{(pdf\_to\_text)} — 2/2}
\small
\begin{tabularx}{\textwidth}{@{}>{\raggedright\arraybackslash}p{3.6cm} >{\raggedright\arraybackslash}X >{\raggedright\arraybackslash}p{4.4cm}@{}}
\toprule
\textbf{규칙} & \textbf{판단 근거} & \textbf{결과 예} \\
\midrule
작은 글씨가 살짝 위에 있으면 위첨자 & 글자 크기·높이 & \texttt{S\&P 500} + \texttt{®} → \texttt{S\&P 500®} \\
줄 간격이 보통이고 문장이 이어지면 같은 문단 & 줄 간격, 앞 줄이 \texttt{,}·\texttt{and}·\texttt{the} 등으로 끝나거나 다음 줄이 소문자로 시작 & 화면의 세 줄 제목 → 한 줄 \\
줄 끝 하이픈 + 소문자로 시작하는 다음 줄 & 하이픈 & \texttt{under-} + \texttt{lying} → \texttt{under-lying} \src{(하이픈은 남김)} \\
\bottomrule
\end{tabularx}
\end{frame}

% ─────────────────────────────────────────────
\begin{frame}{③ 두 종류 공통, 그리고 바꾸지 \alert{않는} 것}
\begin{columns}[T]
\begin{column}{0.47\textwidth}
\begin{block}{③ 두 종류 공통}
\footnotesize
\begin{tabularx}{0.96\linewidth}{@{}X X@{}}
\textbf{규칙} & \textbf{결과} \\
\midrule
날짜 가운데의 줄바꿈은 공백으로 & \texttt{June 29,{\fontspec{Apple Symbols}⏎}2028}\newline → \texttt{June 29, 2028} \\
글머리표({\fontspec{Apple Symbols}●•▪}) 바로 뒤에 글자가 붙으면 공백 하나 & \texttt{{\fontspec{Apple Symbols}●}The notes}\newline → \texttt{{\fontspec{Apple Symbols}●} The notes} \\
\end{tabularx}
\end{block}
\end{column}
\begin{column}{0.47\textwidth}
\begin{alertblock}{바꾸지 않는 것}
\small
\begin{itemize}\setlength\itemsep{3pt}
\item \textbf{글자 자체.} 철자, 숫자, 기호, 원문에 있는 공백\newline \src{(\texttt{5, 210.975}처럼 원문이 이상해도 그대로)}
\item \textbf{``숫자 사이 공백 지우기'' 같은 정리.}\newline 예전 변환(v1)이 이것을 했다가 날짜를 \texttt{October 16,2025}로 망가뜨렸습니다.
\end{itemize}
\end{alertblock}
\end{column}
\end{columns}
\end{frame}

% ─────────────────────────────────────────────
\begin{frame}{변환이 맞는지 어떻게 아나 \src{\normalsize 강사 검사}}
\footnotesize
\renewcommand{\arraystretch}{1.15}
\begin{tabularx}{\textwidth}{@{}l X >{\raggedright\arraybackslash}p{3.3cm}@{}}
\toprule
\textbf{검사} & \textbf{기준} & \textbf{결과} \\
\midrule
1. 글자 보존 & HTML 화면 글자와 변환 결과가 공백·\texttt{\textbar}를 빼면 순서까지 같다 & 19,285건 실패 \num{0} \\
2. 숫자 온전함 & 날짜·금액·비율이 줄바꿈으로 끊기지 않고 원문에 없는 공백이 끼지 않는다 & 19,285건 실패 \num{0} \\
3. 정답 근거 & 정답 27건의 근거 문구 261개가 한 줄 안에 그대로 있다 & \textcolor{accent}{\textbf{261/261}} \\
브라우저 대조 & 크롬으로 그린 글자 위치의 간격과 변환 결과의 띄어쓰기를 비교 \src{(표본 29건)} & 글머리표·넓은 빈칸 문제를 찾아 고침 \\
\bottomrule
\end{tabularx}
\vspace{0.3em}
\begin{alertblock}{알려진 한계}
기호 글꼴로 그린 글머리표는 원문 글자가 영문자라 \texttt{n If the …}처럼 남습니다.
흐름도 안의 칸 몇 곳은 공백으로 이어집니다.
\end{alertblock}
\end{frame}

% ─────────────────────────────────────────────
\begin{frame}[fragile]{내 10월 원문 바꾸기}
\begin{columns}[T]
\begin{column}{0.44\textwidth}
\small Claude Code 패널 또는 VS Code 터미널:
\begin{Verbatim}
uv run jpm_kit/to_text.py
\end{Verbatim}
\vspace{0.5em}
\src{출력에 나오는 예시 accession은 다음 단계에서 씁니다.}
\end{column}
\begin{column}{0.52\textwidth}
\small
\begin{tabularx}{\textwidth}{@{}X p{2.6cm}@{}}
\toprule
\textbf{확인} & \textbf{내 값} \\
\midrule
변환한 공시 수 & \\[6pt] \arrayrulecolor{muted!35}\midrule\arrayrulecolor{black}
일반 HTML / PDF 변환형 & \hfil / \hfil \\[6pt] \arrayrulecolor{muted!35}\midrule\arrayrulecolor{black}
글자 보존 검사 통과 / 실패 & \hfil / \hfil \\[6pt]
\bottomrule
\end{tabularx}
\end{column}
\end{columns}
\end{frame}

% ─────────────────────────────────────────────
\begin{frame}[fragile]{두 종류를 하나씩 원문과 나란히}
\small 위 출력의 예시 accession을 하나씩 씁니다.
\begin{Verbatim}
uv run jpm_kit/to_text.py --show <accession>
\end{Verbatim}
같은 공시를 브라우저로 엽니다: \verb|jpm_kit/sec_oct/html/<accession>.html| 더블클릭.

\footnotesize\renewcommand{\arraystretch}{1.15}
\begin{tabularx}{\textwidth}{@{}X p{3.6cm} p{3.6cm}@{}}
\toprule
 & \textbf{일반 HTML} & \textbf{PDF 변환형} \\
\midrule
accession & & \\[1pt] \arrayrulecolor{muted!35}\midrule\arrayrulecolor{black}
HTML 크기 → 텍스트 크기 & & \\[1pt] \arrayrulecolor{muted!35}\midrule\arrayrulecolor{black}
화면에서 여러 줄인 제목이 텍스트에서는? & & \\[1pt] \arrayrulecolor{muted!35}\midrule\arrayrulecolor{black}
가격 표(Price to Public …) 한 행이 텍스트에서는? & & \\[1pt] \arrayrulecolor{muted!35}\midrule\arrayrulecolor{black}
원문과 다르게 보이는 곳이 있으면 적기 & & \\[1pt]
\bottomrule
\end{tabularx}
\end{frame}

% ─────────────────────────────────────────────
\begin{frame}[fragile]{Claude에게 변환기 한 부분 설명 듣기}
\begin{block}{Claude Code 패널에 붙여 넣기}
\begin{Verbatim}
jpm_kit/htmltext.py 의 _line_text 함수와 pdf_to_text 함수가 하는 일을 설명해줘.
특히 빈 간격 조각을 붙임/공백/칸으로 나누는 기준이 무엇이고, 그 숫자(0.15, 5.0, 1.5)를 바꾸면
어떤 문제가 생길지 예를 들어줘. 코드는 고치지 마.
\end{Verbatim}
\end{block}
\vspace{0.4em}
\begin{alertblock}{확인}
Claude의 설명이 이 실습지 3절의 표와 맞나요? 다르면 코드에서 해당 줄을 찾아 \textbf{누가 맞는지} 확인합니다.
\end{alertblock}
\end{frame}

% ─────────────────────────────────────────────
\begin{frame}{생각해 보기}
\begin{enumerate}\setlength\itemsep{10pt}
\item 태그만 지우는 변환기를 쓰면 다음 실습(규칙으로 값 뽑기)에서 무엇이 틀어질까요?
\item 검사 1은 공백과 \texttt{|}를 빼고 비교합니다. 그래서 검사 1이 \alert{못} 잡는 오류는 어떤 것일까요?
\item 같은 공시를 PDF로 받아 OCR로 읽으면 어떤 점이 더 나쁠까요?
\end{enumerate}
\end{frame}

% 1일차 실습 3 — Azure에 내 DB 만들기 (CLI, 포털, Claude)
% 원본: 실습_Azure_DB_만들기.md · master deck에 \input
\section{실습 3. Azure에 내 DB 만들기}

% ─────────────────────────────────────────────
\begin{frame}{목표: 같은 서버를 세 가지 방법으로}
\begin{block}{목표}
\small 같은 것(\textbf{Azure PostgreSQL 서버})을 세 가지 방법으로 만들어 보고, 방법마다 \textbf{무엇이 보이고 무엇이 남는지} 비교한다.
그다음 내가 바꾼 10월 텍스트를 내 DB에 넣고, 강사의 \textbf{공용 상품 DB}와 이어 본다.
\end{block}
\footnotesize\renewcommand{\arraystretch}{1.1}
\begin{tabularx}{\textwidth}{@{}c X l >{\raggedright\arraybackslash}p{3.6cm}@{}}
\toprule
\textbf{순서} & \textbf{방법} & \textbf{누가} & \textbf{시간} \\
\midrule
\tagbox{accent}{A} & \textbf{CLI}: 명령 한 줄에 설정 전부 & 강사 시연 (VS Code 터미널) & 10분 \\
\tagbox{blue}{B} & \textbf{포털}: 화면에서 칸을 하나씩 골라 만들기 → \alert{내 DB} & 나 & 20분 \src{(만들어지는 6분 동안 C를 본다)} \\
\tagbox{teal}{C} & \textbf{Claude}: 말로 요구사항 → Claude가 명령 스크립트 → 읽고 실행 → 비교 후 삭제 & 강사 시연 (VS Code Claude Code) & 10분 \\
\tagbox{ink}{D} & 내 DB에 10월 텍스트 넣고 VS Code에서 들여다보기 & 나 & 10분 \\
\tagbox{ink}{E} & 공용 상품 DB 읽고 내 공시와 이어 보기 & 나 & 5분 \\
\bottomrule
\end{tabularx}
\end{frame}

% ─────────────────────────────────────────────
\begin{frame}{세 방법, 하나의 결과}
\centering
\begin{tikzpicture}
\node[box=accent,minimum width=4.2cm,minimum height=2.2cm,text width=3.9cm] (a) at (0,0)
  {\tagbox{accent}{A · CLI}\\[3pt] 명령 한 줄에 설정 전부\\[2pt] \scriptsize\color{muted} 강사 · VS Code 터미널 · 10분};
\node[box=blue,minimum width=4.2cm,minimum height=2.2cm,text width=3.9cm] (b) at (4.7,0)
  {\tagbox{blue}{B · 포털}\\[3pt] 화면에서 칸을 하나씩\\[2pt] \scriptsize\color{muted} 나 · portal.azure.com · 20분};
\node[box=teal,minimum width=4.2cm,minimum height=2.2cm,text width=3.9cm] (c) at (9.4,0)
  {\tagbox{teal}{C · Claude}\\[3pt] 말로 요구사항 → 스크립트\\[2pt] \scriptsize\color{muted} 강사 · VS Code Claude Code · 10분};
\node[box=ink,minimum width=8cm,minimum height=1cm] (s) at (4.7,-2.6)
  {\textbf{Azure Database for PostgreSQL} 유연한 서버};
\draw[arr] (a.south) -- (s.north west);
\draw[arr] (b.south) -- (s.north);
\draw[arr] (c.south) -- (s.north east);
\end{tikzpicture}
\vspace{0.6em}

\raggedright\small
\textbf{도구} \src{(사전 준비에서 설치)}: VS Code, VS Code의 \textbf{Claude Code} 확장, \textbf{Azure CLI}(\texttt{az}),
VS Code의 \textbf{PostgreSQL} 확장(Microsoft, \texttt{ms-ossdata.vscode-pgsql}).\\
VS Code 터미널에서 \texttt{az login}을 한 번 해 둔다.
\end{frame}

% ─────────────────────────────────────────────
\begin{frame}[fragile]{A. 강사 시연: CLI로 만들기}
강사가 VS Code 터미널에서 \texttt{tools/create\_pg.sh}를 실행한다.
\alert{명령 하나의 옵션이 곧 설정}이다. 다음 표를 보며 따라 읽고, B에서 포털 칸과 하나씩 맞춰 본다.
\vspace{0.4em}
\begin{Verbatim}
az postgres flexible-server create \
  --resource-group jpm-course --name jpm-pg-demo --location japaneast \
  --tier Burstable --sku-name Standard_B1ms \
  --storage-size 32 --storage-auto-grow Disabled --version 16 \
  --admin-user jpmadmin --admin-password <실행 중에 입력> \
  --public-access <내 공인 IP> --backup-retention 7 --yes
az postgres flexible-server parameter set -g jpm-course -s jpm-pg-demo -n azure.extensions -v VECTOR
\end{Verbatim}
\end{frame}

% ─────────────────────────────────────────────
\begin{frame}{CLI 옵션 ↔ 포털 칸 \src{\normalsize (1/2) 기본 사항}}
\footnotesize
\begin{tabularx}{\textwidth}{@{}>{\raggedright\arraybackslash}p{4.1cm} >{\raggedright\arraybackslash}X >{\raggedright\arraybackslash}p{3.3cm} p{2.3cm}@{}}
\toprule
\textbf{CLI 옵션} & \textbf{뜻} & \textbf{포털에서는} \src{(B에서 채울 칸)} & \textbf{내가 찾은 화면 위치} \\
\midrule
\texttt{-{}-resource-group jpm-course} & 묶음 이름. 지울 때 통째로 & 기본 사항 → 리소스 그룹 & \\[4pt]
\texttt{-{}-name} & 서버 이름 = 주소의 앞부분 (전 세계에서 유일) & 기본 사항 → 서버 이름 & \\[4pt]
\texttt{-{}-location japaneast} & 지역. \alert{한국 지역은 무료 등급이 안 만들어짐} & 기본 사항 → 지역 & \\[4pt]
\texttt{-{}-version 16} & PostgreSQL 버전 & 기본 사항 → 버전 & \\[4pt]
\texttt{-{}-admin-user/}\newline\texttt{-{}-admin-password} & 관리자 계정 & 기본 사항 → 인증 & \\[4pt]
\bottomrule
\end{tabularx}
\end{frame}

% ─────────────────────────────────────────────
\begin{frame}{CLI 옵션 ↔ 포털 칸 \src{\normalsize (2/2) 사양 · 네트워킹 · 만든 뒤}}
\footnotesize\renewcommand{\arraystretch}{1.1}
\begin{tabularx}{\textwidth}{@{}>{\raggedright\arraybackslash}p{4.7cm} >{\raggedright\arraybackslash}X >{\raggedright\arraybackslash}p{3.2cm} p{2.2cm}@{}}
\toprule
\textbf{CLI 옵션} & \textbf{뜻} & \textbf{포털에서는} \src{(B에서 채울 칸)} & \textbf{내가 찾은 화면 위치} \\
\midrule
\texttt{-{}-tier Burstable}\newline\texttt{-{}-sku-name Standard\_B1ms} & 가장 작은 사양 (무료 계정 \num{12개월} 무료) & 컴퓨팅 + 스토리지 → 서버 구성 & \\[1pt]
\texttt{-{}-storage-size 32}\newline\texttt{-{}-storage-auto-grow Disabled} & 저장 공간 \num{32GB}, 자동으로 늘리지 않음 & 컴퓨팅 + 스토리지 & \\[1pt]
\texttt{-{}-backup-retention 7} & 백업 \num{7일} 보관 & 컴퓨팅 + 스토리지 → 백업 & \\[1pt]
\texttt{-{}-public-access <IP>} & 이 IP에서만 접속 허용 (방화벽) & 네트워킹 → 방화벽 규칙 & \\[1pt]
\texttt{parameter set}\newline\texttt{azure.extensions VECTOR} & 벡터 검색 확장(pgvector) 허용 & 만든 뒤 → 서버 매개변수 & \\[1pt]
\bottomrule
\end{tabularx}
\vspace{0.1em}
\small
\tagbox{teal}{CLI의 장점} 같은 명령을 다시 실행하면 같은 서버가 나온다. 기록이 남는다.\\
\tagbox{accent}{단점} 옵션 이름을 알아야 한다.
\end{frame}

% ─────────────────────────────────────────────
\begin{frame}{B. 포털에서 직접 만들기 → 내 DB}
\begin{center}
\begin{tikzpicture}[node distance=0.35cm]
\node[box=blue,minimum width=2.6cm,font=\footnotesize] (p1) {portal.azure.com};
\node[box=blue,minimum width=3.6cm,font=\footnotesize,right=of p1] (p2) {검색창\\\textbf{Azure Database for PostgreSQL}};
\node[box=blue,minimum width=3.2cm,font=\footnotesize,right=of p2] (p3) {\textbf{유연한 서버}\\(Flexible server)};
\node[box=accent,minimum width=1.8cm,font=\footnotesize,right=of p3] (p4) {\textbf{만들기}};
\draw[arr] (p1) -- (p2); \draw[arr] (p2) -- (p3); \draw[arr] (p3) -- (p4);
\end{tikzpicture}
\end{center}
\vspace{0.2em}
\begin{block}{1. 기본 사항}
\begin{itemize}
\item 리소스 그룹: 새로 만들기 → \texttt{jpm-course}
\item 서버 이름: \texttt{jpm-pg-<내 영문 이니셜>-<숫자 2자리>} \src{(예: \texttt{jpm-pg-hgd-07})}. 이미 있으면 숫자를 바꾼다.
\item 지역: \textbf{Japan East} · PostgreSQL 버전: \textbf{16}
\end{itemize}
\end{block}
\end{frame}

% ─────────────────────────────────────────────
\begin{frame}{B-1. 기본 사항 (이어서): 사양 · 인증}
\small
\begin{itemize}
\item 워크로드 유형: \textbf{개발} → \textbf{컴퓨팅 + 스토리지 구성}:\\
  버스트 가능(Burstable), \textbf{Standard\_B1ms}, 저장소 \textbf{32GiB}, 저장소 자동 증가 \textbf{끔}
\end{itemize}
\begin{alertblock}{주의: 기본값은 B2s}
기본값이 \textbf{B2s}(2 vCores, 4 GiB)로 잡힌다. \textbf{서버 구성(Configure server)}을 눌러
\num{Standard\_B1ms (1 vCore, 2 GiB)}로 바꾼다. 무료 계정의 무료 대상은 \textbf{B1ms뿐}이다.
\end{alertblock}
\begin{itemize}
\item 고가용성: 끔
\item 인증: \textbf{PostgreSQL 인증만}, 관리자 이름 \texttt{jpmadmin}, 비밀번호 12자 이상 \alert{(종이에 적기)}
\end{itemize}
\end{frame}

% ─────────────────────────────────────────────
\begin{frame}{B-2. 네트워킹: 내 IP 한 줄만}
\begin{columns}[T]
\begin{column}{0.47\textwidth}
\begin{block}{누른다}
방화벽 규칙:\\
\textbf{+ 현재 클라이언트 IP 주소 추가}\\
\src{(Add current client IP address)}를 \textbf{한 번}.\\[3pt]
아래 표에 \textbf{내 IP 한 줄}이 생기면 된다.
\end{block}
\end{column}
\begin{column}{0.47\textwidth}
\begin{alertblock}{누르지 않는다}
바로 옆의 \textbf{+ 0.0.0.0 - 255.255.255.255 추가}\\
\src{→ 전 세계에 열림}\\[5pt]
\textbf{Azure 서비스의 공용 액세스 허용}은 체크하지 않는다.
\end{alertblock}
\end{column}
\end{columns}
\end{frame}

% ─────────────────────────────────────────────
\begin{frame}{B-3, B-4. 검토 + 만들기 → VECTOR 허용}
\begin{block}{3. 검토 + 만들기}
화면에 나온 \textbf{예상 비용}을 확인한다\\
\src{B1ms면 정가 월 \$23.40, 그 아래에 무료 750시간·32GB 표시}\\[2pt]
→ \textbf{만들기}. \num{약 6분} 걸린다. 기다리는 동안 강사의 C 시연을 본다.
\end{block}
\vspace{0.3em}
\begin{block}{4. 다 만들어지면}
리소스로 이동 → 왼쪽 \textbf{서버 매개변수} → 검색 \texttt{azure.extensions} → \textbf{VECTOR} 체크 → 저장
\end{block}
\end{frame}

% ─────────────────────────────────────────────
\begin{frame}{B. 확인: 내 서버}
\begin{tabularx}{\textwidth}{@{}>{\raggedright\arraybackslash}p{7.6cm} X@{}}
\toprule
\textbf{확인} & \textbf{내 값} \\
\midrule
서버 이름 & \\[8pt]
주소 \src{(개요의 ``서버 이름'' = \texttt{….postgres.database.azure.com})} & \\[8pt]
예상 월 비용 \src{(검토 화면)} & \\[8pt]
만드는 데 걸린 시간 & \\[8pt]
\bottomrule
\end{tabularx}
\end{frame}

% ─────────────────────────────────────────────
\begin{frame}[fragile]{C. 강사 시연: Claude로 만들기 → 비교 → 삭제}
\small 강사가 VS Code의 Claude Code 패널에 \textbf{말로} 요구사항을 주고, Claude가 쓴 스크립트를 함께 읽은 뒤 실행한다. \alert{실행 전에 무엇을 확인하는지} 본다.
\vspace{0.3em}
\begin{Verbatim}
Azure CLI 로 PostgreSQL 유연한 서버를 하나 만드는 스크립트 make_pg2.sh 를 써줘. 실행하지는 마.
조건: 리소스 그룹 jpm-course, 지역 japaneast, 서버 이름 jpm-pg-demo-c,
Burstable Standard_B1ms, 저장소 32GB 자동 증가 끔, PostgreSQL 16, 관리자 jpmadmin,
방화벽은 이 노트북의 공인 IP만(서버를 만들 때 --public-access 에 그 IP를 넣어 한 번에), azure.extensions 에 VECTOR 허용.
비밀번호는 스크립트가 실행될 때 터미널에서 직접 입력하게 하고(화면에 안 보이게), 어디에도 저장하지 마.
명령마다 무엇을 하는지 한 줄 주석을 달아줘.
az 옵션 이름은 기억으로 쓰지 말고, 쓰는 명령마다 이 노트북에서 `az <명령> --help` 를 실행해 확인한 뒤 써. --output 옵션은 붙이지 마.
\end{Verbatim}
\end{frame}

% ─────────────────────────────────────────────
\begin{frame}{C. 왜 그 두 지시가 필요한가}
\small
\begin{alertblock}{괄호 안의 방화벽 지시와 마지막 줄이 없으면}
Claude가 서버만 만들고 \textbf{방화벽 단계에서 멈추는} 일이 잦다.
\end{alertblock}
\vspace{0.2em}
\begin{columns}[T]
\begin{column}{0.47\textwidth}
\begin{block}{방화벽에서 멈추는 이유}
\begin{itemize}
\item 옛 문법(\texttt{firewall-rule create -n <서버>})을 쓴다
\item \texttt{-{}-public-access None}으로 만든 뒤 규칙을 따로 추가하려 한다
\end{itemize}
\src{\texttt{None}은 도움말 설명과 달리 공개 접근이 꺼진 서버를 만들어 규칙을 추가할 수 없다 (az 2.92).}
\end{block}
\end{column}
\begin{column}{0.47\textwidth}
\begin{block}{\texttt{-{}-output table}을 붙이면}
서버 만들기에 붙이면 서버는 생기지만\\
\textbf{``Table output unavailable''} 오류로 스크립트가 멈춰 \alert{VECTOR가 빠진다}.
\end{block}
\end{column}
\end{columns}
\end{frame}

% ─────────────────────────────────────────────
\begin{frame}[fragile]{C. 시연 순서}
\small
\begin{enumerate}
\item Claude가 쓴 스크립트를 \textbf{A의 명령과 한 줄씩 대조}한다. 다른 옵션이 있으면 Claude에게 왜 다른지 묻는다.
\item 강사가 VS Code 터미널에서 실행한다 \src{(비밀번호는 터미널에 직접 입력)}.
\item Claude에게 \textit{``A·C 두 서버의 설정을 az CLI 로 읽어 비교표를 만들고 다른 칸을 표시해줘''} → 비교표를 같이 본다.
\item 두 번째 서버를 지운다:
\end{enumerate}
\begin{Verbatim}
az postgres flexible-server delete -g jpm-course -n jpm-pg-demo-c --yes
\end{Verbatim}
\end{frame}

% ─────────────────────────────────────────────
\begin{frame}{C. 시연을 보며 적기}
\begin{tabularx}{\textwidth}{@{}>{\raggedright\arraybackslash}p{6.4cm} X@{}}
\toprule
\textbf{시연을 보며 적기} & \\
\midrule
Claude 스크립트가 A 명령과 다른 옵션 & \\[14pt]
실행 전에 강사가 확인한 것 & \\[14pt]
세 방법(CLI·포털·Claude)이 각각 남기는 기록 & \\[14pt]
\bottomrule
\end{tabularx}
\end{frame}

% ─────────────────────────────────────────────
\begin{frame}[fragile]{D. 내 DB에 10월 텍스트 넣기}
\small
\textbf{1.} 사전 준비 때 만든 \texttt{jpm\_kit/.env}를 \textbf{열어서} \texttt{PGHOST}(B의 주소)와 \texttt{PGPASSWORD} 두 줄을 \textbf{직접} 채운다.
\begin{alertblock}{주의}
\texttt{.env.example}을 다시 복사하면 \texttt{SEC\_CONTACT}가 지워지니 복사하지 않는다. \textbf{Claude에게 보여 주지 않는다.}
\end{alertblock}
\vspace{0.2em}
\textbf{2.} Claude Code 패널:
\begin{Verbatim}
jpm_kit/load_text.py 가 무엇을 하는지 설명해줘(.env 는 열지 마). 특히 텍스트를 어떤 단위로 잘라 chunks 에 넣는지,
'공통 문구'를 왜 한 번만 넣는지. 설명이 끝나면 uv run jpm_kit/load_text.py 를 실행해줘.
\end{Verbatim}
\end{frame}

% ─────────────────────────────────────────────
\begin{frame}[fragile]{D-3. VS Code에서 내 DB 들여다보기}
\small
VS Code 왼쪽 \textbf{PostgreSQL} 아이콘 → 새 연결 → 내 서버 주소·\texttt{jpmadmin}·비밀번호 → \texttt{postgres} 데이터베이스 → 표 \texttt{filings}, \texttt{chunks}를 열어 본다. 새 쿼리 창에서:
\begin{Verbatim}
select accession, chars, left(body, 120) from filings order by chars limit 5;
select is_boilerplate, count(*), avg(length(body))::int from chunks group by 1;
\end{Verbatim}
\vspace{0.2em}
\begin{tabularx}{\textwidth}{@{}l p{2.6cm} X@{}}
\toprule
\textbf{표} & \textbf{행 수} & \textbf{1행은 무엇인가} \\
\midrule
\texttt{filings} & & \\[6pt]
\texttt{chunks} \src{(공통 문구 수)} & & \\[6pt]
\bottomrule
\end{tabularx}
\end{frame}

% ─────────────────────────────────────────────
\begin{frame}[fragile]{E. 공용 상품 DB 읽기}
\small
상품 표는 강사가 10개월 공시로 만든 것을 \textbf{모두 같이} 쓴다 \src{(읽기 전용)}. 접속 정보는 아래와 같고, 키트의 \texttt{.env.example}에 이미 들어 있다.
\texttt{.env}에 이 네 줄이 있는지 확인한다 \src{(없으면 복사해 넣기)}.
\begin{Verbatim}
SHARED_HOST=jpm-shared-d32c.postgres.database.azure.com
SHARED_USER=student_ro
SHARED_PASSWORD=20261128
SHARED_DATABASE=jpm
\end{Verbatim}
\begin{alertblock}{연세대 네트워크에서만}
읽기 전용 계정이라 표를 바꿀 수 없고, \textbf{연세대 네트워크(교내 와이파이·유선)에서만} 접속된다.
휴대폰 핫스팟이나 집에서는 연결되지 않는다.
\end{alertblock}
\end{frame}

% ─────────────────────────────────────────────
\begin{frame}[fragile]{E. 내 공시와 이어 보기}
그다음:
\begin{Verbatim}
uv run jpm_kit/shared_products.py
\end{Verbatim}
\vspace{0.3em}
\begin{tabularx}{\textwidth}{@{}>{\raggedright\arraybackslash}p{7.6cm} X@{}}
\toprule
\textbf{확인} & \textbf{내 값} \\
\midrule
공용 상품 표의 상품 수 / 공시 수 & \\[8pt]
원금 확인 상품 / 합계 & \\[8pt]
내 10월 공시 중 상품이 있는 공시 / 없는 공시 & \\[8pt]
\bottomrule
\end{tabularx}
\end{frame}

% ─────────────────────────────────────────────
\begin{frame}{생각해 보기}
\begin{columns}[T]
\begin{column}{0.47\textwidth}
\begin{block}{1. 세 방법}
``어제와 똑같은 서버를 내일 다시 만들기''에 가장 좋은 것은? 감사(누가 무엇을 만들었나)에는?
\end{block}
\begin{block}{2. 실행 전 확인}
Claude가 만든 스크립트를 실행하기 전에 꼭 확인해야 할 것 두 가지는? \src{(C 시연에서 본 것)}
\end{block}
\end{column}
\begin{column}{0.47\textwidth}
\begin{block}{3. 빈 공시}
내 10월 공시인데 상품 표에 상품이 없는 공시는 무엇일까?
\end{block}
\begin{block}{4. 읽기 전용}
공용 상품 DB가 읽기 전용인 이유는?
\end{block}
\end{column}
\end{columns}
\end{frame}

% 1일차 실습 4 — 임베딩 → SQL로 묻기 → 더 복잡한 질의 (실습_SQL_임베딩_질의.md)
\section{실습 4. 임베딩 → SQL로 묻기}
% 글꼴에 ⓪(U+24EA)이 없어 TikZ로 대신 그린다
\providecommand{\circzero}{\tikz[baseline=(z.base)]\node[circle,draw,line width=0.4pt,inner sep=0.4pt,minimum size=0.95em,font=\scriptsize](z){0};}

% ─────────────────────────────────────────────
\begin{frame}{임베딩 → SQL로 묻기 → 더 복잡한 질의 \src{\normalsize 1일차 마무리}}
\begin{block}{목표}
글을 숫자로 바꾸는 \textbf{임베딩}이 무엇인지 손으로 확인하고, JPMorgan 상품 담당자가 되어 상품 라인업을 \textbf{SQL로} 분석한 뒤, \textbf{벡터로만} 할 수 있는 분석(특이 상품 찾기, 새 상품과 닮은 상품 찾기)을 해 본다.
\end{block}
\vspace{0.3em}
\footnotesize
\begin{tabularx}{\textwidth}{@{}l X r@{}}
\toprule
\textbf{순서} & \textbf{내용} & \textbf{시간} \\
\midrule
\tagbox{accent}{A} & Azure OpenAI 만들고 \textbf{벡터화 시작} (약 19분 걸리므로 먼저 시작해 두고 B·C 동안 돌린다) & 8분 \\
\tagbox{blue}{B} & 임베딩이 무엇인가: TensorFlow Embedding Projector로 보기, 단어 계산, 문장 유사도 & 15분 \\
\tagbox{teal}{C} & 상품군 분석: JPMorgan 상품 담당자로서 ``무엇을 얼마나 팔았나'' (SQL, 벡터 불필요) & 12분 \\
\tagbox{ink!55}{D} & 벡터로만 할 수 있는 분석: 키워드 vs 뜻(3분), 특이 상품 찾기, 새 상품 설계안과 닮은 상품 찾기 & 18분 \\
\tagbox{muted}{E} & 1일차 정리 & 5분 \\
\bottomrule
\end{tabularx}
\end{frame}

% ═════════════════════════════════════════════ A
\begin{frame}[fragile]{\tagbox{accent}{A} Azure OpenAI 만들기}
VS Code 터미널 \src{(리소스 그룹 \texttt{jpm-course}, 이름은 \texttt{jpm-aoai-<이니셜>-<숫자>})}
\vspace{0.4em}
\begin{Verbatim}
az cognitiveservices account create -g jpm-course -n jpm-aoai-<이니셜>-<숫자> -l japaneast \
  --kind OpenAI --sku S0 --custom-domain jpm-aoai-<이니셜>-<숫자> --yes
az cognitiveservices account deployment create -g jpm-course -n jpm-aoai-<이니셜>-<숫자> \
  --deployment-name emb3s --model-name text-embedding-3-small --model-version 1 --model-format OpenAI \
  --sku-name GlobalStandard --sku-capacity 1000
\end{Verbatim}
\vspace{0.6em}
\begin{columns}[T]
\begin{column}{0.48\textwidth}
\begin{block}{만드는 것}
\small 계정 하나 + 임베딩 모델 배포 \texttt{emb3s}\\
\src{모델 \texttt{text-embedding-3-small}, 지역 \texttt{japaneast}}
\end{block}
\end{column}
\begin{column}{0.48\textwidth}
\begin{alertblock}{먼저 시작해 둔다}
\small 벡터화는 \num{약 19분}. B·C 동안 돌린다.
\end{alertblock}
\end{column}
\end{columns}
\end{frame}

\begin{frame}[fragile]{\tagbox{accent}{A} \texttt{InsufficientQuota} 오류가 나면}
\begin{alertblock}{무료 계정은 Azure OpenAI를 \textbf{하나만}}
지운 계정도 \num{48시간} 동안 ``삭제 대기''로 남아 그 자리를 차지한다\\
\src{(예: 연습으로 만들었다 지운 계정)}
\end{alertblock}
\vspace{0.4em}
휴지통을 확인하고 완전히 지운 뒤 앞의 명령을 다시 실행한다.
\vspace{0.4em}
\begin{Verbatim}
az cognitiveservices account list-deleted -o table
az cognitiveservices account purge -l japaneast -g <그 계정의 리소스 그룹> -n <그 계정 이름>
\end{Verbatim}
\end{frame}

\begin{frame}[fragile]{\tagbox{accent}{A} 키 적고 벡터화 시작}
\small
\texttt{.env}에 \textbf{직접} 적는다.
\begin{itemize}\setlength\itemsep{1pt}
\item \texttt{AOAI\_ENDPOINT=https://jpm-aoai-<이니셜>-<숫자>.openai.azure.com/}
\item \texttt{AOAI\_KEY} \src{(포털 → 이 리소스 → \textbf{키 및 엔드포인트} → 키 1)}
\end{itemize}
\vspace{0.3em}
\begin{Verbatim}
uv run jpm_kit/embed.py --try 3
\end{Verbatim}
청크 3개가 각각 \num{숫자 1,536개}로 바뀌는 것을 본다. 그다음 Claude Code 패널:
\vspace{0.3em}
\begin{Verbatim}
jpm_kit/embed.py 를 백그라운드로 실행하고, jpm_kit/out/embed_progress.txt 로 진행 상황을 알려줘.
\end{Verbatim}
\end{frame}

% ═════════════════════════════════════════════ B
\begin{frame}{\tagbox{blue}{B} 임베딩이 무엇인가 \src{\normalsize 벡터화가 도는 동안}}
\begin{block}{한 줄 정의}
글(단어·문장·문단)을 \textbf{숫자 목록(벡터)}으로 바꾸는 것. 뜻이 비슷한 글은 비슷한 숫자 목록이 되도록 학습되어 있어서, \textbf{숫자끼리의 거리로 뜻이 얼마나 가까운지} 잴 수 있다.
\end{block}
\vspace{0.5em}
\centering
\begin{tikzpicture}
\node[box=ink,minimum width=3.2cm] (t1) {글 하나};
\node[box=ink,minimum width=3.2cm,below=0.5cm of t1] (t2) {뜻이 비슷한 글};
\node[box=blue,minimum width=5.6cm,right=1.8cm of t1,font=\small\ttfamily] (v1) {[ 숫자, 숫자, 숫자, … ]};
\node[box=blue,minimum width=5.6cm,right=1.8cm of t2,font=\small\ttfamily] (v2) {[ 숫자, 숫자, 숫자, … ]};
\draw[arr] (t1) -- node[above,font=\scriptsize]{임베딩} (v1);
\draw[arr] (t2) -- node[above,font=\scriptsize]{임베딩} (v2);
\draw[accent,line width=0.9pt,{Stealth[length=5pt]}-{Stealth[length=5pt]}] ([xshift=2mm]v1.east) to[bend left=50] node[right,font=\scriptsize\bfseries,accent,align=left]{거리가\\가깝다} ([xshift=2mm]v2.east);
\end{tikzpicture}
\end{frame}

\begin{frame}{\tagbox{blue}{B} 말 네 개}
\begin{tabularx}{\textwidth}{@{}l X@{}}
\toprule
\textbf{말} & \textbf{뜻} \\
\midrule
벡터 & 숫자 목록. 단어 벡터는 \num{200개}, 오늘 쓰는 문장 임베딩은 \num{1,536개} \\
차원 & 숫자의 개수. 숫자 하나하나에는 사람이 읽을 뜻이 없다 \\
코사인 유사도 & 두 벡터의 방향이 얼마나 같은가. \textbf{1}이면 같은 방향(뜻이 가까움), \textbf{0}이면 무관 \\
벡터 검색 & 질문도 같은 모델로 벡터로 바꾼 뒤, 가장 가까운 청크를 찾는 것 \\
\bottomrule
\end{tabularx}
\end{frame}

\begin{frame}{\tagbox{blue}{B} \circzero{} 눈으로 보기: TensorFlow Embedding Projector \src{\normalsize 브라우저, 5분}}
\small
구글 TensorFlow 팀이 공개한 임베딩 시각화 도구. 위키백과 글로 학습한 단어 \textbf{1만 개}의 벡터(단어마다 숫자 200개)를 3차원으로 눌러 점으로 보여 준다.
\begin{enumerate}\setlength\itemsep{3pt}
\item \texttt{projector.tensorflow.org}를 연다. 왼쪽 위 데이터가 \textbf{Word2Vec 10K}인지 확인한다(기본값).
\item 오른쪽 \textbf{Search} 칸에 \texttt{bank} → 아래 목록에서 \textbf{bank를 클릭} → 오른쪽에 \textbf{Nearest points in the original space}(숫자 200개 공간에서 가장 가까운 단어)와 거리가 나온다.\\
\src{이 숫자는 \textbf{코사인 거리}라 \textbf{작을수록 가깝다} (① 스크립트가 보여 주는 \textbf{유사도}는 클수록 가까움).}
\item \textbf{Isolate 101 points}를 눌러 이웃만 남기고, 점 구름을 마우스로 돌려 본다. \texttt{dollar}, \texttt{bond}, \texttt{risk}도 찾아본다.
\item 왼쪽 아래 투영 방법을 \textbf{PCA → T-SNE}로 바꿔, 비슷한 단어끼리 무리 짓는 모습을 본다.
\end{enumerate}
\end{frame}

\begin{frame}{\tagbox{blue}{B} \circzero{} 적을 것}
\begin{tabularx}{\textwidth}{@{}X p{6cm}@{}}
\toprule
\textbf{확인} & \textbf{내 값} \\
\midrule
\texttt{bank}의 가장 가까운 단어 5개 & \\[4pt]
\texttt{bank} 이웃 중 금융이 아닌 단어 (어떤 뜻의 bank?) & \\[4pt]
\texttt{bond}의 가장 가까운 단어 3개 (채권 뜻인가?) & \\[4pt]
\bottomrule
\end{tabularx}
\vspace{0.6em}
\begin{alertblock}{단어 하나에 숫자 묶음이 하나뿐}
\small \texttt{bank} 이웃에는 banks·monetary·imf 같은 금융 단어 사이에 river·thames·canal(강둑), gaza·west·strip(요르단강 서안 West Bank)이 섞여 나온다. 뜻이 섞인 것이다 → ②에서 문장 임베딩이 이것을 어떻게 푸는지 본다.
\end{alertblock}
\end{frame}

\begin{frame}[fragile]{\tagbox{blue}{B} ① 같은 자료로 계산하기: 단어 계산}
\small
AI 호출 없음, 내 노트북에서 numpy로. \circzero{}에서 본 \textbf{바로 그 Word2Vec 10K 자료}(TensorFlow Embedding Projector 공개 파일, 약 8MB)를 내려받아 직접 계산한다. 핵심은 아래 몇 줄.
\vspace{0.3em}
\begin{Verbatim}
idx = {word: i for i, word in enumerate(words)}          # 단어 → 행 번호
q = V[idx["king"]] - V[idx["man"]] + V[idx["woman"]]     # 벡터 산술
scores = (V @ q) / (np.linalg.norm(V, axis=1) * np.linalg.norm(q))   # 모든 단어와 코사인 유사도
# scores 가 가장 큰 단어 = 답 (식에 쓴 세 단어는 뺀다)
\end{Verbatim}
\vspace{0.3em}
\begin{Verbatim}
uv run jpm_kit/emb_1_words.py
\end{Verbatim}
\end{frame}

\begin{frame}{\tagbox{blue}{B} ① 무엇이 나오나}
\begin{columns}[T]
\begin{column}{0.48\textwidth}
\begin{block}{가까운 단어}
\texttt{dollar → currency, usd, franc …}\\
\src{\circzero{}의 목록과 비교해 보기}
\end{block}
\begin{block}{단어 계산}
\texttt{king - man + woman = queen}\\
\texttt{paris - france + japan = tokyo}\\[2pt]
→ 관계가 \textbf{방향}으로 들어 있다
\end{block}
\end{column}
\begin{column}{0.48\textwidth}
\begin{alertblock}{한계}
\texttt{bond}의 이웃이 hydrogen, molecule(화학 결합).\\[2pt]
\texttt{bank}는 loan(\num{0.50})보다 river(\num{0.56}, 강둑)에 가깝다\\[2pt]
→ 단어 하나에 숫자 묶음이 하나뿐이라 뜻이 섞인다
\end{alertblock}
\src{\texttt{jpm\_kit/out/word\_map.png}: 200개 숫자를 2차원으로 눌러 본 지도}
\end{column}
\end{columns}
\end{frame}

\begin{frame}[fragile]{\tagbox{blue}{B} ① 왕·여왕·남자·여자 그림 \src{\normalsize \texttt{jpm\_kit/out/king\_queen.png}}}
\begin{columns}[c]
\begin{column}{0.44\textwidth}
\centering
\begin{tikzpicture}[x=1cm,y=1cm]
\draw[muted,-{Stealth[length=4pt]}] (0,0) -- (4.4,0) node[below left,font=\scriptsize,align=right]{남자→여자\\(woman - man)};
\draw[muted,-{Stealth[length=4pt]}] (0,0) -- (0,3.2) node[above,font=\scriptsize]{보통 사람→왕족};
\node[circle,fill=ink,inner sep=1.5pt,label={[font=\scriptsize]below:man}] (m) at (0.9,0.7) {};
\node[circle,fill=ink,inner sep=1.5pt,label={[font=\scriptsize]below:woman}] (w) at (3.3,0.7) {};
\node[circle,fill=accent,inner sep=1.5pt,label={[font=\scriptsize]above:king}] (k) at (0.9,2.4) {};
\node[circle,fill=accent,inner sep=1.5pt,label={[font=\scriptsize]above:queen}] (q) at (3.3,2.4) {};
\draw[arr,ink] (m) -- (w);
\draw[arr,accent] (k) -- (q);
\end{tikzpicture}\\
\src{화살표 방향만 나타낸 모식도}
\end{column}
\begin{column}{0.52\textwidth}
\small
가로축 = 남자→여자 방향(woman - man), 세로축 = 보통 사람→왕족 방향.\\[3pt]
man→woman 화살표와 \textbf{king→queen 화살표가 같은 쪽(오른쪽)}을 향한다. 그래서 king에서 man을 빼고 woman을 더하면 queen 근처에 닿는다.\\[3pt]
prince→princess, boy→girl, father→mother도 같은 쪽.\\[3pt]
\alert{같은 쪽을 향하지 않는 쌍(husband→wife)}도 찾아 왜 그럴지 생각해 본다.
\end{column}
\end{columns}
\vspace{0.4em}
\src{Projector에서 같은 그림을 보려면 (강사 시연): 왼쪽 아래 \textbf{CUSTOM} → 네 칸 오른쪽의 \texttt{.*} 버튼을 눌러 켠 뒤 Left \texttt{\^{}man\$}, Right \texttt{\^{}woman\$}, Up \texttt{\^{}(king|queen)\$}, Down \texttt{\^{}(man|woman)\$}}
\end{frame}

\begin{frame}[fragile]{\tagbox{blue}{B} ① 직접 해 보기}
\small
두 명령을 \textbf{따로} 실행한다. \alert{한 줄에 같이 쓰면 \texttt{-{}-near}만 실행된다.}
\vspace{0.3em}
\begin{Verbatim}
uv run jpm_kit/emb_1_words.py --calc tokyo japan france
uv run jpm_kit/emb_1_words.py --near loan
\end{Verbatim}
\vspace{0.4em}
\begin{tabularx}{\textwidth}{@{}l X@{}}
\toprule
\texttt{-{}-calc A B C} & A - B + C에 가장 가까운 단어 \src{(tokyo - japan + france → paris)} \\
\texttt{-{}-near 단어} & 그 단어와 가장 가까운 단어 6개 \src{(loan → loans, debts, payments …)} \\
\bottomrule
\end{tabularx}

\vspace{0.6em}
다른 영어 단어로도 해 본다.
\end{frame}

\begin{frame}[fragile]{\tagbox{blue}{B} ② 문장을 숫자로 \src{\normalsize 내 Azure OpenAI, 1원도 안 듦}}
\begin{Verbatim}
uv run jpm_kit/emb_2_sentences.py
\end{Verbatim}
\vspace{0.4em}
\begin{tabularx}{\textwidth}{@{}X p{5.5cm}@{}}
\toprule
\textbf{확인} & \textbf{내 값} \\
\midrule
같은 bond: 채권 문장끼리 vs 채권–화학 문장 & \hfill / \hfill\null \\[6pt]
한국어 질문과 가장 가까운 영어 공시 문장 & \\[4pt]
금리만 다른 두 문장(10.75\% vs 8.10\%)의 유사도 & \\[4pt]
\bottomrule
\end{tabularx}

\vspace{0.5em}
\small 내 문장으로:
\begin{Verbatim}
uv run jpm_kit/emb_2_sentences.py "원금 보장" "principal protected" "capital at risk"
\end{Verbatim}
\end{frame}

\begin{frame}{\tagbox{blue}{B} ② 문장 임베딩이 보여 주는 것}
\begin{columns}[T]
\begin{column}{0.31\textwidth}
\begin{block}{문맥}
\textbf{문맥}이 뜻을 정한다\\
\src{단어 벡터의 한계를 넘음}
\end{block}
\end{column}
\begin{column}{0.31\textwidth}
\begin{block}{언어}
\textbf{언어가 달라도} 뜻이 같으면 가깝다
\end{block}
\end{column}
\begin{column}{0.31\textwidth}
\begin{alertblock}{숫자}
대신 \textbf{정확한 숫자는 거의 구별하지 못한다}
\end{alertblock}
\end{column}
\end{columns}
\vspace{1em}
\begin{center}
→ 그래서 \num{키워드 검색과 합친다}.
\end{center}
\vspace{0.6em}
\src{벡터화는 계속 돌고 있다. C(SQL)를 마칠 무렵 끝난다. 결과는 D 전에 적는다.}
\end{frame}

% ═════════════════════════════════════════════ C
\begin{frame}{\tagbox{teal}{C} 상품군 분석: 우리는 무엇을 얼마나 팔았나 \src{\normalsize 벡터화가 도는 동안}}
\begin{block}{역할}
여러분은 \textbf{JPMorgan 구조화상품 담당자}. 공용 상품 DB(\texttt{products}, 2026년 1–10월 \num{8,868개})로 우리 상품 라인업을 본다.\\[3pt]
상품군 이름을 이미 알고 있으므로 \textbf{벡터가 필요 없다. SQL이면 된다.}
\end{block}
\vspace{0.5em}
\textbf{공용 상품 DB 연결} \src{— PostgreSQL 확장에 두 번째 연결 (교내 네트워크에서만)}\\[4pt]
\small
\begin{tabular}{@{}l l@{}}
\toprule
서버 & \texttt{jpm-shared-d32c.postgres.database.azure.com} \\
사용자 & \texttt{student\_ro} \\
비밀번호 & \texttt{20261128} \\
데이터베이스 & \texttt{jpm} \\
\bottomrule
\end{tabular}
\end{frame}

\begin{frame}{\tagbox{teal}{C-1} 내 손으로 SQL 두 개 \src{\normalsize VS Code 쿼리 창, 5분}}
\footnotesize
\begin{tabularx}{\textwidth}{@{}c >{\raggedright\arraybackslash}p{4.2cm} >{\raggedright\arraybackslash}X p{2.6cm}@{}}
\toprule
\textbf{\#} & \textbf{질문} & \textbf{힌트} & \textbf{내 답} \\
\midrule
1 & 상품군(\texttt{product\_type})별 상품 수와 비중 & \texttt{count(*)}, \texttt{round(100.0 * count(*) / sum(count(*)) over (), 1)}, \texttt{group by 1 order by 2 desc} & \\[4pt]
2 & 조건부 쿠폰 상품의 기초자산 구조(\texttt{underlier\_structure})별 평균 연 쿠폰 — \textbf{분모(쿠폰이 있는 상품 수)도 함께} & \texttt{where product\_type = 'contingent\_income' and coupon\_annual\_pct is not null group by …} & \\[4pt]
\bottomrule
\end{tabularx}
\vspace{0.8em}
\begin{block}{막히면}
\small Claude Code 패널에 \textit{``이 질문을 SQL로 어떻게 쓰는지 설명만 해줘''}라고 묻고, 받은 SQL을 \textbf{읽고 이해한 뒤} 실행한다.
\end{block}
\end{frame}

\begin{frame}[fragile]{\tagbox{teal}{C-2} Claude에게 상품군 지도 맡기기 \src{\normalsize Claude Code 패널, 약 4분}}
\begin{Verbatim}[fontsize=\footnotesize]
너는 JPMorgan 구조화상품 담당자다. 공용 상품 DB(products, 2026년 1~10월)로 "우리가 무엇을 얼마나 팔았나" 상품군 지도를 만들어줘.
1. 상품군(product_type)별 상품 수·비중·평균 만기·평균 연 쿠폰·발행가 대비 추정가치 차이(price_value_gap_pct) 평균·원금 합계. 평균마다 분모(값이 있는 상품 수)도.
2. 월별 상품 수와 상위 4개 상품군의 월별 추이.
3. 가장 많이 쓰인 기초자산(tickers) 10개.
4. 이 숫자로 말할 수 있는 것 3가지와, 이 데이터로는 말할 수 없는 것 2가지.
DB 연결은 jpm_kit/kit_db.py 의 connect("SHARED"). 짧은 코드는 jpm_kit/out/ 에 파일로 만들어 uv run 으로 실행하고, 쓴 SQL 을 모두 보여줘. .env 는 열지 마.
\end{Verbatim}
\end{frame}

\begin{frame}{\tagbox{teal}{C-2} 확인}
\begin{tabularx}{\textwidth}{@{}X p{5.6cm}@{}}
\toprule
\textbf{확인} & \\
\midrule
Claude의 상품군별 상품 수가 내 C-1의 1번과 같은가 & \\[8pt]
평균 쿠폰·숨은 비용(gap)의 \textbf{분모}가 상품 수와 다른 상품군은? 왜 다른가 & \\[8pt]
Claude가 꼽은 ``이 데이터로는 말할 수 없는 것'' 하나 & \\[8pt]
\bottomrule
\end{tabularx}
\end{frame}

% ═════════════════════════════════════════════ D
\begin{frame}{\tagbox{ink!55}{D} 벡터로만 할 수 있는 분석 \src{\normalsize 벡터화가 끝난 뒤}}
먼저 벡터화 결과를 적는다 \src{(\texttt{jpm\_kit/out/embed\_progress.txt} 마지막 줄)}
\vspace{0.6em}
\begin{tabularx}{\textwidth}{@{}X p{7cm}@{}}
\toprule
\textbf{벡터화 결과} & \textbf{내 값} \\
\midrule
청크 수 / 걸린 시간 / 토큰 / 비용 & \\[8pt]
DB 크기 (전 → 후) & \\[8pt]
\bottomrule
\end{tabularx}
\end{frame}

\begin{frame}{\tagbox{ink!55}{D} 키워드 검색 vs 벡터 검색}
\small
\begin{tabularx}{\textwidth}{@{}l X X@{}}
\toprule
 & \textbf{키워드 검색} & \textbf{벡터 검색} \\
\midrule
찾는 기준 & 질문과 \textbf{같은 단어}가 들어 있는가 \src{(\texttt{tsv} 키워드 색인)} & 질문과 \textbf{뜻}이 가까운가 \src{(\texttt{emb} 벡터의 코사인 거리)} \\
잘 찾는 것 & 공시 용어 그대로, 정확한 숫자, 고유명사 & 바꿔 말한 질문 (공시와 다른 말로 물어도) \\
놓치는 것 & 질문의 단어가 한 문단에 모두 있어야 찾음 & 숫자만 다른 문장을 구별 못 함 (B의 10.75\% vs 8.10\%), 한국어로 막연히 물으면 자주 실패 \\
\bottomrule
\end{tabularx}
\end{frame}

\begin{frame}{\tagbox{ink!55}{D} 하이브리드: 두 결과의 \textbf{순위}를 합친다}
\begin{columns}[c]
\begin{column}{0.56\textwidth}
\centering
\begin{tikzpicture}
\node[box=accent,minimum width=3cm] (k) at (0,0.85) {키워드 검색 결과\\[-1pt]\scriptsize 1위, 2위, 3위, …};
\node[box=blue,minimum width=3cm] (v) at (0,-0.85) {벡터 검색 결과\\[-1pt]\scriptsize 1위, 2위, 3위, …};
\node[box=ink,minimum width=2.6cm,minimum height=2.2cm] (h) at (4.6,0) {하이브리드\\[-1pt]\scriptsize 점수 합으로\\[-3pt]\scriptsize 다시 줄 세움};
\draw[arr] (k.east) -- node[above,sloped,font=\scriptsize]{$1/(60+r)$} (h.west);
\draw[arr] (v.east) -- node[below,sloped,font=\scriptsize]{$1/(60+r)$} (h.west);
\end{tikzpicture}
\end{column}
\begin{column}{0.40\textwidth}
\begin{block}{RRF}
\small \texttt{ask.py}의 세 번째 결과.\\[2pt]
각 목록에서 \textbf{r위}면 \num{1 / (60 + r)} 점.
\end{block}
\begin{alertblock}{늘 이기지는 않는다}
\small 양쪽에 모두 오른 문단이 위로 올라오지만, 늘 이기지는 않는다.
\end{alertblock}
\end{column}
\end{columns}
\end{frame}

\begin{frame}[fragile]{\tagbox{ink!55}{D-1} 같은 뜻, 다른 말: \texttt{ask.py}로 3분}
\begin{Verbatim}
uv run jpm_kit/ask.py "Can the bank end the investment early?" -k 3
uv run jpm_kit/ask.py "6.0% per annum daily deduction" -k 3
\end{Verbatim}
\vspace{0.6em}
\begin{tabularx}{\textwidth}{@{}l X X@{}}
\toprule
\textbf{질문} & \textbf{키워드 1위} & \textbf{벡터 1위} \\
\midrule
고객의 말 (공시에 없는 표현) & & \\[8pt]
공시 표현 그대로 & & \\[8pt]
\bottomrule
\end{tabularx}
\end{frame}

\begin{frame}{\tagbox{ink!55}{D-1} 첫 질문에서 키워드가 아무것도 못 찾는 이유}
\begin{block}{키워드 검색은 질문의 \textbf{모든 단어}가 한 문단에 있어야 한다}
\texttt{'bank' \& 'end' \& 'invest' \& 'earli'}\\[3pt]
벡터는 \textbf{``Early redemption risk''} 문단을 찾는다.
\end{block}
\vspace{0.6em}
\begin{alertblock}{벡터가 값을 하는 때}
\textbf{묻는 말과 공시의 말이 다를 때}.\\[3pt]
아래 두 분석(D-2, D-3)이 바로 그런 경우: 무엇을 찾을지 \textbf{이름을 모르는} 분석.
\end{alertblock}
\end{frame}

\begin{frame}{\tagbox{ink!55}{D} 공시 한 건 = 벡터 하나 \src{\normalsize \texttt{jpm\_kit/similar.py}}}
\centering
\begin{tikzpicture}
\node[box=ink,minimum width=2.4cm,minimum height=2.6cm] (doc) {공시 1건};
\node[box=muted,minimum width=2.6cm,minimum height=0.55cm,right=1.2cm of doc,yshift=0.85cm,font=\scriptsize] (c1) {청크 벡터};
\node[box=muted,minimum width=2.6cm,minimum height=0.55cm,right=1.2cm of doc,font=\scriptsize] (c2) {청크 벡터};
\node[box=muted,minimum width=2.6cm,minimum height=0.55cm,right=1.2cm of doc,yshift=-0.85cm,font=\scriptsize] (c3) {청크 벡터 …};
\node[box=blue,minimum width=2.4cm,right=1.2cm of c2] (avg) {평균\\[-1pt]\scriptsize = 벡터 하나};
\draw[arr] (doc.east) -- (c1.west);
\draw[arr] (doc.east) -- (c2.west);
\draw[arr] (doc.east) -- (c3.west);
\draw[arr] (c1.east) -- (avg.west);
\draw[arr] (c2.east) -- (avg.west);
\draw[arr] (c3.east) -- (avg.west);
\node[font=\scriptsize,muted,below=0.15cm of c3] {공통 문구가 아닌 청크만};
\node[box=accent,minimum width=3cm,right=0.8cm of avg,font=\small] (cmp) {코사인 유사도로\\[-1pt]비교};
\draw[arr] (avg) -- (cmp);
\end{tikzpicture}
\vspace{0.8em}
\begin{itemize}\setlength\itemsep{2pt}\small
\item 공시 \num{1,859건}을 벡터 \num{1,859개}로 만든다 \src{(처음 한 번 약 30초)}
\item 공시끼리, 또는 새 설명과 코사인 유사도로 비교한다
\end{itemize}
\end{frame}

\begin{frame}[fragile]{\tagbox{ink!55}{D-2} 우리가 판 상품 중 특이한 것이 있나 \src{\normalsize 약 5분}}
먼저 직접 실행해 본다.
\begin{Verbatim}
uv run jpm_kit/similar.py --outliers 5
\end{Verbatim}
그다음 Claude Code 패널에서:
\begin{Verbatim}[fontsize=\footnotesize]
너는 JPMorgan 구조화상품 담당자다. 10월에 우리가 낸 공시 중 특이한 상품이 있는지 보려고 한다.
1. uv run jpm_kit/similar.py --outliers 5 를 실행하고 결과를 표로 보여줘.
2. 그 5건이 어떤 상품인지 uv run jpm_kit/ask.py "<검색어>" --accession <번호> 로 원문을 찾아 확인하고, 왜 특이한지 한 줄씩 써줘.
   문서 양식(판매 채널)이 달라서인지, 상품 구조가 달라서인지 구분하고, 근거 [accession #chunk_id] 를 붙여.
.env 는 열지 마.
\end{Verbatim}
\end{frame}

\begin{frame}{\tagbox{ink!55}{D-2} 확인}
\begin{tabularx}{\textwidth}{@{}X p{5.6cm}@{}}
\toprule
\textbf{확인} & \\
\midrule
가장 닮은 다른 공시와 0.99 이상으로 닮은 공시 비율 → 무슨 뜻인가 & \\[8pt]
특이 5건 중 \textbf{상품 구조}가 특이한 것 / \textbf{문서 양식(판매 채널)}이 달라서인 것 & \hfill / \hfill\null \\[8pt]
특이 5건 중 공용 상품 표에 없거나 분류 미확정인 것은 몇 건? 왜 그럴까 & \\[8pt]
\bottomrule
\end{tabularx}
\end{frame}

\begin{frame}[fragile]{\tagbox{ink!55}{D-3} 새 상품을 설계했다: 기존 상품 중 무엇과 닮았나 \src{\normalsize 약 5분}}
\small Claude Code 패널에서:
\begin{Verbatim}[fontsize=\footnotesize]
너는 JPMorgan 구조화상품 담당자다. 새 상품을 설계했다:
"금과 비트코인 ETF를 같은 비중으로 담은 바스켓에 연동, 만기 5년, 원금 100% 보장, 바스켓 상승분의 80% 참여, 수익 상한 40%, 이자 없음."
1. 이 설명을 공시에서 쓰는 영어 문장 3~4개로 바꿔 보여줘.
2. 비교 기준으로 평범한 상품 설명(Nasdaq-100·Russell 2000 worst-of 자동조기상환 조건부 쿠폰)도 영어로 하나 써서,
   uv run jpm_kit/similar.py --text "<새 상품>" --text "<평범한 상품>" 을 실행해.
3. 새 상품과 가장 비슷한 공시 3건이 무엇 때문에 닮았고 무엇이 다른지 ask.py --accession 으로 원문을 확인해 표로. 특히 "원금 100% 보장"이 닮은 이유에 들어 있는지 봐.
4. 이 새 상품은 기존 상품의 변형인가, 새로운 구조인가? 근거와 함께 3줄로.
.env 는 열지 마.
\end{Verbatim}
\end{frame}

\begin{frame}{\tagbox{ink!55}{D-3} 확인}
\begin{tabularx}{\textwidth}{@{}X p{5.6cm}@{}}
\toprule
\textbf{확인} & \\
\midrule
새 상품 / 평범한 상품의 1위 유사도 & \hfill / \hfill\null \\[8pt]
새 상품과 가장 닮은 공시가 닮은 이유 (무엇이 같은가) & \\[8pt]
``원금 100\% 보장''이 닮은 이유에 들어갔나? 들어가지 않았다면 그것이 벡터에 대해 말해 주는 것은 & \\[8pt]
\bottomrule
\end{tabularx}
\vspace{0.6em}
\begin{alertblock}{한국어 설명은 먼저 영어 공시 문장으로}
\small 한국어 설명을 그대로 \texttt{-{}-text}에 넣으면 영어 공시와 잘 비교되지 않는다 \src{(리허설: 유사도 \num{0.35})}. 그래서 Claude에게 먼저 영어 공시 문장으로 바꾸게 한다.
\end{alertblock}
\end{frame}

% ═════════════════════════════════════════════ E
\begin{frame}{\tagbox{muted}{E} 1일차 정리 \src{\normalsize 5줄 메모}}
\begin{enumerate}\setlength\itemsep{7pt}
\item 오늘 내 손으로 만든 것: SEC 원문 \underline{\hspace{1.2cm}}건 → 텍스트 \underline{\hspace{1.2cm}}건 → 내 DB의 표 \underline{\hspace{1.2cm}}개 (행 수)
\item 변환기가 바꾸지 \textbf{않는} 것 한 가지
\item 키워드 검색이 잘 찾는 질문 / 벡터 검색이 잘 찾는 질문
\item 벡터가 놓친 것 한 가지와 그 이유 \src{(예: D-3의 원금 보장, B의 숫자)}
\item 2일차에 확인하고 싶은 질문
\end{enumerate}
\vspace{0.6em}
\begin{alertblock}{2일차에는 Azure를 쓰지 않는다 — 수업 후 지워도 된다}
VS Code 터미널에서 \texttt{az group delete -n jpm-course -{}-yes} \src{(내 DB와 Azure OpenAI가 함께 지워진다)}
\end{alertblock}
\end{frame}


% ─────────────────────────────────────────────
\section{정리}

\begin{frame}{오늘 한 일}
\centering
\begin{tikzpicture}[node distance=0.5cm]
\node[box=accent,minimum width=3.1cm,text width=2.9cm] (a) {\textbf{받기}\\\scriptsize Claude와 만든 코드로\\SEC 10월 원문};
\node[box=blue,minimum width=3.1cm,text width=2.9cm,right=of a] (b) {\textbf{바꾸기}\\\scriptsize 글자는 그대로,\\사이만 정한다};
\node[box=teal,minimum width=3.1cm,text width=2.9cm,right=of b] (c) {\textbf{넣기}\\\scriptsize CLI · 포털 · Claude\\세 가지로 만든 DB};
\node[box=ink,minimum width=3.1cm,text width=2.9cm,right=of c] (d) {\textbf{묻기}\\\scriptsize SQL로 상품군,\\벡터로 닮음};
\draw[arr] (a) -- (b); \draw[arr] (b) -- (c); \draw[arr] (c) -- (d);
\end{tikzpicture}
\vspace{1.4em}

\begin{columns}[T]
\begin{column}{0.44\textwidth}
\begin{block}{정리 (5분)}
실습 4의 E — 1일차 정리 메모를 쓴다.
\end{block}
\end{column}
\begin{column}{0.44\textwidth}
\begin{alertblock}{퀴즈 (15분)}
오픈북 · 객관식 5 + 서술형 2
\end{alertblock}
\end{column}
\end{columns}
\end{frame}

\end{document}
